% !TeX root = ../../Book2.tex
\chapter{振荡积分}
\label{b2:ch:21}
\BookChapterNumber{b2:ch:21}
\begin{chapterabstract}
振荡积分的大小取决于哪些部分能够相互抵消。相位没有驻点时，分部积分把抵消转为衰减，但边界和振幅导数必须一并控制；相位有驻点时，局部尺度、Hesse 矩阵与振幅消失阶决定主要贡献。本章从这些标量问题出发，再用 \(TT^*\) 方法处理允许输入自由振荡的算子，区分普通 Hesse 矩阵、混合 Hesse 矩阵和相位驻点方程的三种作用。
\end{chapterabstract}

\begin{learninggoals}\begin{itemize}
\item\label{b2:item:21-goal-1}区分端点项、驻点贡献和非驻相余项，说明支撑与振幅导数如何进入估计常数。
\item\label{b2:item:21-goal-2}由次水平集和分部积分理解 van der Corput 衰减，辨别有限型、单调性与参数尺度的作用。
\item\label{b2:item:21-goal-3}通过 Morse 坐标和复 Gaussian 计算驻相系数，核对行列式、符号、Jacobi 因子及参数余项。
\item\label{b2:item:21-goal-4}用混合导数下界、\(TT^*\) 与 Schur 估计证明局部 \(L^2\) 衰减，并判断局部化和参数一致性的适用范围。
\end{itemize}\end{learninggoals}

积分 \(\int e^{2\pi\mathrm i\lambda\phi}\cdot a\) 中的指数因子始终具有模一。对它直接取绝对值，只能得到振幅的总量，无法看到 \(\lambda\) 增大后的抵消。一个可行的起点是观察相位的一阶导数：它不为零时，可以把指数写成一个导数；它等于零时，则应先找出附近的自然尺度。由此产生非驻相法与驻相法，但两种方法都需要知道振幅在哪里非零。

\ProseParagraphBreak

本章先在一维把边界项完整保留下来，再证明 van der Corput 引理。随后由复 Gaussian 与光滑换元得到多维驻相展开，并讨论临界点碰撞、振幅随尺度变化和曲面曲率。一般上界与渐近首项承担不同任务：一个统一上界未必给出非零主项，而几个各自非零的驻点贡献也可能相互抵消。

\ProseParagraphBreak

进入\secref{b2:sec:2104}后，对象从一个固定振幅的积分变成作用于整个 \(L^2\) 单位球的算子。输入可以选择自己的相位，不能把标量驻相公式直接用于每一个输入。\(TT^*\) 把问题改写成两个输出点之间的相位差，混合 Hesse 矩阵由此出现。该矩阵在一个点可逆，只先保证足够小的邻域内有统一估计；把小片相加与直接在整个支撑上使用同一个梯度下界，是不同的论证。

\ProseParagraphBreak

分部积分从相位导数中提取衰减，Morse 换元则把非退化驻点附近的相位化成二次型；两种方法分别处理远离驻点和接近驻点的区域。选读的\secref{b2:sec:2105}进一步用相位记录位置与频率之间的对应关系，介绍 Fourier 积分算子。前七道习题围绕主线安排；最后一道用实际波包检验半波传播的符号与尺度。

\ProseParagraphBreak

Elias M. Stein 的《Harmonic Analysis: Real-Variable Methods, Orthogonality, and Oscillatory Integrals》\cite{Stein1993Harmonic}可用于继续阅读振荡积分及其在调和分析中的应用。本章保留指数中的 \(2\pi\)，使用其他归一化的资料时，应同时换算驻相系数和传播参数。

% !TeX root = ../../Book2.tex
\section{一维振荡积分}
\label{b2:sec:2101}

正积分的大小由振幅总量决定，振荡积分还受到相位变化的制约。对有限区间 \([\alpha,\beta]\)、实值相位 \(\phi\) 和复值振幅 \(a\)，记
\[
 I_\lambda(a,\phi)
 \coloneqq\int_\alpha^\beta
 e^{2\pi\mathrm i\lambda\phi(t)}\cdot a(t)\,\mathrm dt,
 \quad \lambda>0.
\]
本章主要考察 \(\lambda\) 很大的情形；本节的一阶和高阶导数估计对任意正 \(\lambda\) 都成立。这一点在把导数下界吸收到参数后尤其有用，因为新的有效参数未必大于一。

\ProseParagraphBreak

直接取绝对值只能得到 \(|I_\lambda|\leq\|a\|_1\)，完全看不到振荡。要获得衰减，须比较邻近位置的相位：导数远离零时，相邻波动可以通过分部积分抵消；导数接近零的区域则不能这样处理。van der Corput 方法把这两个区域分开，并用高阶导数控制后一部分的长度。

\subsection{分部积分}
\label{b2:subsec:210101}

先在足够光滑的情形中辨认边界项。设 \(\phi\in C^2([\alpha,\beta];\mathbb R)\)、\(|\phi'|\geq c>0\)，并设 \(a:[\alpha,\beta]\to\mathbb C\) 绝对连续。此时 \(a/\phi'\) 也绝对连续，由
\[
 e^{2\pi\mathrm i\lambda\phi(t)}
 =\frac1{2\pi\mathrm i\lambda\phi'(t)}
   \frac{\mathrm d}{\mathrm dt}e^{2\pi\mathrm i\lambda\phi(t)}
\]
得到
\[
 \begin{aligned}
 I_\lambda(a,\phi)
 &=\left[
 \frac{e^{2\pi\mathrm i\lambda\phi(t)}\cdot a(t)}
      {2\pi\mathrm i\lambda\phi'(t)}
 \right]_\alpha^\beta\\
 &\hspace{1em}
 -\frac1{2\pi\mathrm i\lambda}
 \int_\alpha^\beta e^{2\pi\mathrm i\lambda\phi(t)}
       \cdot\left(\frac a{\phi'}\right)'(t)\,\mathrm dt.
 \end{aligned}
\]
商的导数是 \(a'/\phi'-a\phi''/(\phi')^2\)，所以
\[
 |I_\lambda(a,\phi)|
 \leq\frac1{2\pi\lambda}
 \left[
 \frac{|a(\alpha)|+|a(\beta)|+\|a'\|_1}{c}
 +\frac1{c^2}\int_\alpha^\beta|a(t)|\cdot|\phi''(t)|\,\mathrm dt
 \right].
\]
这个公式清楚地区分了振幅变化、相位导数变化和两个端点的贡献。相位本身很大并不提供下界；给 \(\phi\) 加上一个任意大的常数，只会改变整体复相位。

\ProseParagraphBreak

振幅在端点消失时，边界项为零；否则不能将其省略。例如
\[
 \int_0^1e^{2\pi\mathrm i\lambda t}\,\mathrm dt
 =\frac{e^{2\pi\mathrm i\lambda}-1}{2\pi\mathrm i\lambda}.
\]
沿 \(\lambda=n+1/2\) 的序列，其模为 \(1/(\pi\lambda)\)。这里相位没有驻点，振幅甚至是常数，但仍只有一阶衰减。将一个光滑函数硬截断在闭区间上，与使振幅光滑地消失在端点附近，是两种不同的操作。

\subsection{非驻相}
\label{b2:subsec:210102}

若振幅有紧支撑，导数转移不会产生边界项，可以反复使用刚才的恒等式。为便于后面处理参数族，下面把估计常数所需的振幅范数也写出来。

\Needspace{6\baselineskip}
\begin{proposition}[非驻相衰减]\label{b2:prop:21-nonstationary-one}
设 \(\alpha<\beta\)，\(a\in C_c^\infty((\alpha,\beta);\mathbb C)\)，\(\phi\in C^\infty((\alpha,\beta);\mathbb R)\)。设 \(V\) 是 \(\operatorname{supp}a\) 的一个邻域，在 \(V\) 上有 \(|\phi'|\geq c>0\)。对任意整数 \(N\geq1\)、任意 \(\lambda>0\)，
\[
 \left|\int_\alpha^\beta e^{2\pi\mathrm i\lambda\phi(t)}
                \cdot a(t)\,\mathrm dt\right|
 \leq C_N\lambda^{-N}\sum_{j=0}^N\|a^{(j)}\|_1.
\]
若在 \(V\) 上的 \(\phi',\ldots,\phi^{(N+1)}\) 绝对值不超过 \(M\)，则 \(C_N\) 只依赖 \(N,c^{-1},M\)，不依赖 \(\lambda\) 或振幅。所有 \(L^1\) 范数均在 \((\alpha,\beta)\) 上取。
\end{proposition}
\begin{proof}
令 \(a_0=a\)，并在 \(V\) 上递归定义
\[
 a_{j+1}=-\left(\frac{a_j}{\phi'}\right)'.
\]
因 \(a_j\) 的支撑始终包含于 \(\operatorname{supp}a\)，可将这些函数在 \(V\) 外延为零，得到光滑紧支撑函数。每次积分分部都没有边界项，因此
\[
 I_\lambda(a,\phi)
 =(2\pi\mathrm i\lambda)^{-N}
   \int_\alpha^\beta e^{2\pi\mathrm i\lambda\phi(t)}
          \cdot a_N(t)\,\mathrm dt.
\]
递归中的每次求导至多增加一阶振幅导数；系数则由 \(\phi'\) 的倒数与相位的各阶导数相乘、相加得到。归纳可写成
\[
 a_N=\sum_{j=0}^N c_{N,j}\cdot a^{(j)},
\]
其中每个 \(c_{N,j}\) 只含至多 \(N+1\) 阶的相位导数及 \((\phi')^{-1}\) 的有限次幂。题设给这些系数一致上界，故
\(\|a_N\|_1\leq C\sum_{j=0}^N\|a^{(j)}\|_1\)。
取绝对值即得结论。
\end{proof}

这里“非驻相”指相位一阶导数在振幅支撑附近有正下界。相位在别处可以有驻点，因为那里没有振幅。反之，一个参数族中的驻点即使尚未进入支撑，只要到支撑的距离趋于零，相位导数下界就可能变小，估计常数也会相应增大。

\subsection{van der Corput 引理}
\label{b2:subsec:210103}

高阶导数下界允许一阶导数消失，因此不能在整个区间上使用倒导数。先建立只需要 \(C^1\) 正则性的一阶估计。单调导数未必经典可微，直接写出 \(\phi''\) 并积分并不合法；改用相位本身作变量，可以避开这一问题。

\Needspace{6\baselineskip}
\begin{lemma}[一阶导数估计]\label{b2:lem:21-monotone-first-derivative}
设 \(\alpha<\beta\)，\(\phi\in C^1([\alpha,\beta];\mathbb R)\)，\(|\phi'|\geq c>0\)。若连续函数 \(1/\phi'\) 在该区间的总变差为有限数 \(V\)，则对任意 \(\lambda>0\) 与子区间 \([u,v]\subset[\alpha,\beta]\)，
\[
 \left|\int_u^v e^{2\pi\mathrm i\lambda\phi(t)}\,\mathrm dt\right|
 \leq\frac{2c^{-1}+V}{2\pi\lambda}.
\]
这里总变差是所有有限分割上
\(\sum_j|q(t_j)-q(t_{j-1})|\) 的上确界，其中 \(q=1/\phi'\)。
特别地，若 \(\phi'\) 单调，则上述积分不超过 \(3/(2\pi\lambda c)\)。
\end{lemma}
\begin{proof}
导数连续且不为零，所以它符号恒定，\(\phi\) 严格单调。相位换元 \(s=\phi(t)\) 将积分写成
\[
 \int_{\phi(u)}^{\phi(v)}
       e^{2\pi\mathrm i\lambda s}\cdot q_*(s)\,\mathrm ds,
 \quad
 q_*(s)=\frac1{\phi'(\phi^{-1}(s))}.
\]
若上下限倒置，先反转积分方向，绝对值不变。函数 \(q_*\) 连续，\(\|q_*\|_\infty\leq c^{-1}\)，且沿单调换元的总变差不超过 \(V\)。

在相位区间上取网格大小趋于零的分割，用 \(q_*\) 在分割点的值作分段线性插值 \(q_n\)。连续性保证 \(q_n\) 一致收敛到 \(q_*\)；插值不增大上确界范数，而且
\[
 \int|q_n'(s)|\,\mathrm ds
 =\sum_j|q_*(s_j)-q_*(s_{j-1})|\leq V.
\]
每个 \(q_n\) 绝对连续，普通积分分部给出
\[
 \left|\int e^{2\pi\mathrm i\lambda s}
                 \cdot q_n(s)\,\mathrm ds\right|
 \leq\frac{2\|q_n\|_\infty+\|q_n'\|_1}{2\pi\lambda}
 \leq\frac{2c^{-1}+V}{2\pi\lambda}.
\]
在固定的有限相位区间上，积分随一致逼近收敛，得到 \(q_*\) 的结论。若 \(\phi'\) 单调，倒数也单调，并且符号固定、绝对值至多为 \(c^{-1}\)，故其总变差不超过 \(c^{-1}\)。
\end{proof}

第二个准备事实说明：一个高阶导数始终不小的函数，不可能在很长的集合上保持很小。所控制的是集合的总长度，不要求这个集合本身是一个区间。

\Needspace{6\baselineskip}
\begin{lemma}[次水平集的长度]\label{b2:lem:21-sublevel-measure}
设 \(r\geq1\) 为整数，\(\alpha<\beta\)，\(h\in C^r([\alpha,\beta];\mathbb R)\)，且 \(|h^{(r)}|\geq c>0\)。对任意 \(\delta>0\)，有
\[
 \left|\{\,t\in[\alpha,\beta]:|h(t)|\leq\delta\,\}\right|
 \leq 2r\left(\frac{\delta}{c}\right)^{1/r}.
\]
\end{lemma}
\begin{proof}
记左侧集合为 \(E\)，其测度为 \(L\)。\(L=0\) 时无需证明。设 \(L>0\)，则 \(E\) 是非空紧集。函数
\[
 F(t)=|E\cap[\alpha,t]|
\]
连续且为 \(1\)-Lipschitz 函数，从 \(0\) 增至 \(L\)。取 \(t_0=\min E\)、\(t_r=\max E\)。对 \(1\leq j\leq r-1\)，令 \(t_j\) 为满足 \(F(t)\geq jL/r\) 的最小点。连续性给出 \(F(t_j)=jL/r\)；该点必在 \(E\) 中，否则其邻域不交 \(E\)，\(F\) 在此邻域为常数，与最小性矛盾。端点也分别满足 \(F(t_0)=0\)、\(F(t_r)=L\)。因此，对 \(0\leq i<j\leq r\)，
\[
 t_j-t_i\geq F(t_j)-F(t_i)=\frac{(j-i)L}{r}.
\]
这些点既落在 \(E\) 中，又具有可计算的间距。

用这 \(r+1\) 个点作 Lagrange 插值多项式
\[
 P(t)=\sum_{j=0}^r h(t_j)
       \prod_{\substack{0\leq\ell\leq r\\\ell\ne j}}
       \frac{t-t_\ell}{t_j-t_\ell}.
\]
它的 \(t^r\) 系数为
\[
 A=\sum_{j=0}^r
       \frac{h(t_j)}
       {\prod_{\ell\ne j}(t_j-t_\ell)}.
\]
因为 \(|h(t_j)|\leq\delta\)，且
\[
 \left|\prod_{\ell\ne j}(t_j-t_\ell)\right|
 \geq (L/r)^r\,j!(r-j)!,
\]
所以
\[
 |A|
 \leq\delta\left(\frac rL\right)^r
       \sum_{j=0}^r\frac1{j!(r-j)!}
 =\frac{\delta}{r!}\left(\frac{2r}{L}\right)^r.
\]
函数 \(h-P\) 在 \(t_0,\ldots,t_r\) 为零。先在每两个相邻零点之间应用 Rolle 定理，得到至少 \(r\) 个一阶导数零点；如此重复 \(r\) 次，得到某个 \(\theta\in(t_0,t_r)\)，使
\[
 0=(h-P)^{(r)}(\theta)
   =h^{(r)}(\theta)-r!A.
\]
于是 \(c\leq r!|A|\leq\delta(2r/L)^r\)，整理即得所述长度界。
\end{proof}

现在可以证明\term[van-der-corput-yinli]{van der Corput 引理}[van der Corput lemma]。一阶情形需要控制倒导数的变化；二阶及以上情形则由最高阶导数下界自动限制低阶导数反复转向的次数。

\Needspace{6\baselineskip}
\begin{theorem}[van der Corput 引理]\label{b2:thm:21-vdc}
设 \(\alpha<\beta\)，\(k\geq1\) 为整数，\(c>0\)、\(\lambda>0\)，\(\phi\in C^k([\alpha,\beta];\mathbb R)\)，\(a:[\alpha,\beta]\to\mathbb C\) 绝对连续。若 \(k=1\)，假设 \(|\phi'|\geq c\) 且 \(\phi'\) 单调；若 \(k\geq2\)，假设 \(|\phi^{(k)}|\geq c\)。则
\[
 \left|\int_\alpha^\beta
 e^{2\pi\mathrm i\lambda\phi(t)}\cdot a(t)\,\mathrm dt\right|
 \leq C_k(\lambda c)^{-1/k}
       \bigl(\|a\|_\infty+\|a'\|_1\bigr),
\]
其中 \(C_k\) 只依赖 \(k\)，不依赖区间长度、相位、振幅、\(\lambda\) 或 \(c\)。
\end{theorem}
\begin{proof}
先取振幅恒为一。\(k=1\) 时直接用\cref{b2:lem:21-monotone-first-derivative}。以下设 \(k\geq2\)，取任意阈值 \(\delta>0\)，将区间分为
\[
 E_\delta=\{\,t:|\phi'(t)|\leq\delta\,\}
 \quad\text{和}\quad
 [\alpha,\beta]\setminus E_\delta.
\]
对 \(h=\phi'\)、\(r=k-1\) 应用\cref{b2:lem:21-sublevel-measure}，得到
\[
 \left|\int_{E_\delta}e^{2\pi\mathrm i\lambda\phi(t)}
                         \,\mathrm dt\right|
 \leq |E_\delta|
 \leq 2(k-1)\left(\frac{\delta}{c}\right)^{1/(k-1)}.
\]

补集上的积分使用一阶估计，但须先控制其单调区间数。若 \(k=2\)，\(\phi''\) 连续且不为零，因而符号固定，\(\phi'\) 在整个区间单调。若 \(k>2\)，\(\phi''\) 至多有 \(k-2\) 个互异零点：否则从 \(k-1\) 个零点出发，连续应用 \(k-2\) 次 Rolle 定理，会产生一个 \(\phi^{(k)}\) 的零点，矛盾。因此可将区间分成至多 \(k-1\) 个 \(\phi'\) 单调的区间。在每一片内，\(\phi'>\delta\) 和 \(\phi'<-\delta\) 各至多占一个区间。补集于是由至多 \(2(k-1)\) 个区间组成；端点的开闭不影响积分。在每个这样的区间上，先对内部闭子区间使用一阶引理，再令端点趋向边界，得到
\[
 \left|\int_{[\alpha,\beta]\setminus E_\delta}
             e^{2\pi\mathrm i\lambda\phi(t)}\,\mathrm dt\right|
 \leq \frac{C_k}{\lambda\delta}.
\]
两部分合并为
\[
 \left|\int_\alpha^\beta
          e^{2\pi\mathrm i\lambda\phi(t)}\,\mathrm dt\right|
 \leq C_k\left[
       \left(\frac{\delta}{c}\right)^{1/(k-1)}
       +\frac1{\lambda\delta}\right].
\]
选
\(\delta=c^{1/k}\lambda^{-(k-1)/k}\)，两项都等于
\((\lambda c)^{-1/k}\) 的常数倍。因此该估计适用于全部正的 \(\lambda,c\)，包括有效参数 \(\lambda c<1\) 的情形。

最后处理振幅。上述结论在任意子区间上仍成立，因为导数假设在限制后不变。令
\[
 F(t)=\int_\alpha^t e^{2\pi\mathrm i\lambda\phi(s)}\,\mathrm ds.
\]
则 \(\|F\|_\infty\leq C_k(\lambda c)^{-1/k}\)。函数 \(a\) 绝对连续，积分分部给出
\[
 I_\lambda(a,\phi)
 =a(\beta)F(\beta)-\int_\alpha^\beta a'(t)F(t)\,\mathrm dt,
\]
从而得到所述上界。
\end{proof}

一阶中的单调性可以替换为受控的总变差，但不能仅保留 \(|\phi'|\geq c\)。沿用一阶引理，并作同样的振幅积分分部，可得
\[
 |I_\lambda(a,\phi)|
 \leq \frac{2c^{-1}+\operatorname{Var}(1/\phi')}{2\pi\lambda}
       \bigl(\|a\|_\infty+\|a'\|_1\bigr).
\]
因此关键不是导数是否有零点这一件事，还包括倒导数能否反复变化而积累很大的总变差。

\ProseParagraphBreak

下面的光滑相位族给出真实的失败例子。对正整数 \(N\)，令
\[
 q_N(s)=\frac{2+\cos(2\pi Ns)}3,
 \quad
 T_N(s)=\int_0^s q_N(v)\,\mathrm dv,
 \quad 0\leq s\leq1.
\]
有 \(1/3\leq q_N\leq1\)、\(T_N(0)=0\)、\(T_N(1)=2/3\)，所以 \(T_N\) 有光滑反函数
\(\phi_N:[0,2/3]\to[0,1]\)，并且
\[
 1\leq\phi_N'(t)
 =\frac1{q_N(\phi_N(t))}\leq3.
\]
然而在固定区间上取常数振幅、取 \(\lambda=N\)，换元 \(t=T_N(s)\) 给出
\[
 \begin{aligned}
 \int_0^{2/3}e^{2\pi\mathrm iN\phi_N(t)}\,\mathrm dt
 &=\int_0^1e^{2\pi\mathrm iNs}\cdot
                  \frac{2+\cos(2\pi Ns)}3\,\mathrm ds\\
 &=\frac16.
 \end{aligned}
\]
最后一步把余弦写成两个复指数；所有非零整数频率的全周期积分为零，只留下常数项。故不存在只依赖导数下界的统一 \(C/N\) 估计。此例的倒导数总变差恰为 \(4N/3\)，也解释了为什么变差版本没有产生矛盾。

\subsection{衰减估计}
\label{b2:subsec:210104}

导数下界 \(c\) 与高频参数始终以有效参数 \(\lambda c\) 一起出现。即使 \(\lambda\geq1\)，当相位接近退化时也可能有 \(\lambda c<1\)；上一定理仍然成立，只是此时其右侧可能不如平凡界 \(\|a\|_1\)。实际应用可取两者的较小值。不能先只证明参数大于一的定理，再不加说明地将它用于 \(\lambda c<1\)。

\ProseParagraphBreak

相位 \(t^k\) 展示了幂次的尺度来源。换元 \(t=\lambda^{-1/k}u\) 后，
\[
 \int e^{2\pi\mathrm i\lambda t^k}\cdot a(t)\,\mathrm dt
 =\lambda^{-1/k}
   \int e^{2\pi\mathrm i u^k}\cdot a(\lambda^{-1/k}u)\,\mathrm du.
\]
缩放确定了观察驻点的自然邻域。要取得非零渐近首项，还需要控制扩张积分区间上的振荡尾部，并排除对称性抵消。二次情形的非零主项将在\secref{b2:sec:2102}中算出，从而把自然尺度与实际衰减阶数联系起来。

\ProseParagraphBreak

有限阶导数条件也可局部使用。若 \(\phi\) 在紧支撑附近光滑，且每一点至少有某个 \(1\leq j\leq k\) 满足 \(|\phi^{(j)}|\geq c>0\)，则连续性给出有限个邻域，使每一片上某个固定阶导数有下界 \(c/2\)。乘上光滑分割后，\(j\geq2\) 的片用 van der Corput 引理；\(j=1\) 的片用紧支撑非驻相估计，无须额外假设该片导数单调。对 \(\lambda\geq1\)，有限项相加得到 \(O(\lambda^{-1/k})\)。覆盖、分割函数导数、相位的有限阶导数与下界都参与常数，不能仅记录一个抽象的“有限型”条件。

\ProseParagraphBreak

这种统一上界在驻点碰撞时比逐点二次展开更稳定。取
\[
 I(\lambda,s)=\int_{\mathbb R}
 e^{2\pi\mathrm i\lambda(t^3-st)}\cdot a(t)\,\mathrm dt,
 \quad a\in C_c^\infty(\mathbb R;\mathbb C).
\]
三阶导数恒为六，因此 van der Corput 引理给出
\(|I(\lambda,s)|\leq C_a\lambda^{-1/3}\)，其中 \(C_a\) 与 \(s\in\mathbb R\) 无关。若 \(s>0\)，两个驻点是 \(\pm\sqrt{s/3}\)，各点的二阶导数大小为 \(2\sqrt{3s}\)。后面的二次驻相公式因而会出现 \(s^{-1/4}\lambda^{-1/2}\) 的系数，不能在 \(s\downarrow0\) 时保持统一。

\ProseParagraphBreak

换元 \(t=\lambda^{-1/3}u\) 后，指数中的相位成为
\(u^3-(s\lambda^{2/3})u\)。因此碰撞区域的自然参数是 \(s\lambda^{2/3}\)，相应范围是 \(s\asymp\lambda^{-2/3}\)。在此范围，驻点间距 \(\sqrt{s}\) 与三次振荡尺度 \(\lambda^{-1/3}\) 相当，两个局部二次模型已不能统一分离。对于固定的小正 \(s\)，二次驻点的局部宽度约为 \(\lambda^{-1/2}s^{-1/4}\)；要求它远小于 \(\sqrt{s}\)，也恰好得到 \(s\lambda^{2/3}\gg1\)。

\ProseParagraphBreak

\cref{b2:fig:21-critical-collision}画的是相位导数，而不是原相位；其横轴交点直接显示驻点的数目和碰撞。统一三阶估计无需写出 Airy 函数，尺度比较已经说明了分别使用二次近似的范围。

\begin{figure}[htbp]
\centering
\begin{tikzpicture}[x=1cm,y=.6cm]
\foreach \shift/\offset/\paneltitle in {-3.2/.65/{\(s<0\)},0/0/{\(s=0\)},3.2/-.65/{\(s>0\)}}{
\begin{scope}[shift={(\shift,0)}]
\draw[->] (-1.25,0)--(1.35,0) node[right]{\(t\)};
\draw[AnalysisBlue,thick,domain=-1.05:1.05,samples=45] plot(\x,{1.3*\x*\x+\offset});
\node at (0,2.35){\paneltitle};
\end{scope}}
\fill (0,0) circle(1.3pt);
\fill (2.493,0) circle(1.3pt);
\fill (3.907,0) circle(1.3pt);
\node[below,font=\small] at (2.493,-.95){\(-\sqrt{s/3}\)};
\node[below,font=\small] at (3.907,-.95){\(\sqrt{s/3}\)};
\node at (0,-2.7){碰撞尺度：\(\sqrt{s}\asymp\lambda^{-1/3}\)，即 \(s\lambda^{2/3}\asymp1\)};
\end{tikzpicture}
\caption[三次相位的驻点碰撞]{\(\phi_s'(t)=3t^2-s\) 的三个符号情形，不按比例。横轴交点是驻点；当两个交点的间距降到三次振荡尺度时，不能再把它们作为两个统一非退化的二次模型。}
\label{b2:fig:21-critical-collision}
\end{figure}

非驻相区域仍需连同支撑一起判断。若 \(s<0\)，则
\(\phi_s'(t)=3t^2+|s|\geq|s|\)，固定负参数给出任意阶衰减，但其常数随 \(|s|\downarrow0\) 增长。若振幅支撑满足 \(|t|\geq\varepsilon>0\)，则对足够小的 \(|s|\)，一阶导数仍有统一正下界，可以恢复对这些参数一致的快速衰减。退化的影响由相位、支撑与参数范围共同决定。

% !TeX root = ../../Book2.tex
\section{驻相法}
\label{b2:sec:2102}

一阶导数消失后，非驻相分部积分的分母发生奇异，主要贡献便可能集中在驻点附近。若驻点 \(x_0\) 的二次项非退化，位移 \(h\) 使相位变化约为 \(\lambda h^{\mathsf T}Hh/2\)；在各特征值固定且非零时，变化为常数量级的位移尺度是 \(\lambda^{-1/2}\)。这解释了每个非退化方向为何贡献半个衰减幂次。本节先把远离驻点的部分除去，再用可逆换元将驻点附近的相位精确化为二次型，最后计算其常数与余项。

\subsection{非退化临界点}
\label{b2:subsec:210201}

\Needspace{6\baselineskip}
\begin{definition}\label{b2:def:21-critical}
设整数 \(d\geq1\)，\(U\subset\mathbb R^d\) 为开集，\(\phi\in C^2(U;\mathbb R)\)，\(x_0\in U\)。若 \(\nabla\phi(x_0)=0\)，称 \(x_0\) 为驻点；若同时
\(\det D^2\phi(x_0)\ne0\)，称其为\term[fei-tuihua-zhudian]{非退化驻点}[nondegenerate critical point]。
\end{definition}
逆函数定理用于映射 \(\nabla\phi\)，说明非退化驻点在一个小邻域内是唯一的。反过来，孤立驻点未必非退化：\(t^4\) 的原点就是例子。第二种情形的衰减阶由更高阶项决定，不能套用二次驻相公式。

\ProseParagraphBreak
设相位在紧集 \(K\) 的一个开邻域内光滑。若 \(K\) 中所有驻点非退化，则其中驻点只有有限个。否则一列互异驻点在 \(K\) 内有聚点，梯度的连续性使聚点仍为驻点，与局部孤立性矛盾。于是可以分别计算各驻点贡献，再对剩余部分使用非驻相法；不同贡献的复相位可能抵消，因此各项上界之和只是上界。

\ProseParagraphBreak

先把后一步所需的多维工具在此证明。不能直接在多个坐标中同时除以各个偏导数，因为梯度非零并不保证每一个偏导数都非零。沿梯度方向微分，则只需除以整个梯度的平方长度。

\Needspace{6\baselineskip}
\begin{proposition}[多维非驻相衰减]\label{b2:prop:21-nonstationary-multi}
设整数 \(d\geq1\)，\(U\subset\mathbb R^d\) 开，\(a\in C_c^\infty(U;\mathbb C)\)，\(\phi\in C^\infty(U;\mathbb R)\)。设在 \(K=\operatorname{supp}a\) 的一个邻域 \(W\Subset U\) 内有 \(|\nabla\phi|\geq c>0\)。对每个整数 \(N\geq1\) 及 \(\lambda\geq1\)，
\[
 \left|\int_U e^{2\pi\mathrm i\lambda\phi(x)}
                   \cdot a(x)\,\mathrm dx\right|
 \leq C_N\lambda^{-N}
       \sum_{|\alpha|\leq N}\|\partial^\alpha a\|_{L^1(U)}.
\]
常数 \(C_N\) 只依赖 \(d,N,c^{-1}\) 及 \(\phi\) 在 \(W\) 上至多 \(N+1\) 阶导数的上界。因而固定支撑、相同梯度下界与这些导数的一致界给出一致的估计。
\end{proposition}
\begin{proof}
在 \(W\) 内令
\[
 V(x)=\frac{\nabla\phi(x)}{|\nabla\phi(x)|^2},
 \quad
 Ta=-\operatorname{div}(Va).
\]
因为 \(V\cdot\nabla\phi=1\)，有
\[
 V\cdot\nabla e^{2\pi\mathrm i\lambda\phi}
 =2\pi\mathrm i\lambda e^{2\pi\mathrm i\lambda\phi}.
\]
振幅紧支于 \(W\)，散度积分没有边界项，故
\[
 \int e^{2\pi\mathrm i\lambda\phi}\cdot a
 =(2\pi\mathrm i\lambda)^{-1}
   \int e^{2\pi\mathrm i\lambda\phi}\cdot Ta.
\]
这里使用双线性积分分部，\(T\) 是 \(V\cdot\nabla\) 的形式转置；没有对振幅取复共轭。具体地，
\[
 Ta=-V\cdot\nabla a-(\operatorname{div}V)a.
\]
每个 \(T^ja\) 的支撑仍包含于 \(K\)，故可以重复 \(N\) 次，得到
\[
 \int e^{2\pi\mathrm i\lambda\phi}\cdot a
 =(2\pi\mathrm i\lambda)^{-N}
   \int e^{2\pi\mathrm i\lambda\phi}\cdot T^Na.
\]
用乘积公式归纳，\(T^Na\) 是 \(\partial^\alpha a\)、\(|\alpha|\leq N\)，乘上 \(V\) 的至多 \(N\) 阶导数所组成的有限和。对 \(V=\nabla\phi/|\nabla\phi|^2\) 求导，分母只有 \(|\nabla\phi|\) 的负幂，分子只用到 \(\phi\) 的至多 \(N+1\) 阶导数。题设因此给出这些系数的统一上界，取各项的 \(L^1\) 范数即得结论。
\end{proof}

在驻点附近选一个等于一的光滑截断后，剩余振幅仍紧支撑；若其支撑上没有驻点，紧性便给出梯度的正下界。这正是下面把远离驻点的部分吸收到任意阶余项中的依据。对有边界的积分域，必须另行保留边界项，本命题不适用于直接的示性函数截断。

\subsection{Morse 归约与相位二次型}
\label{b2:subsec:210202}

对实对称可逆矩阵 \(H\)，记其正特征值数减负特征值数为 \(\operatorname{sgn}H\)，称为这里使用的符号差。Taylor 展开只说明相位近似于一个二次型；\term[morse-yinli]{Morse 引理}[Morse lemma]给出更强的局部事实：通过改变坐标，可以把所有高阶项吸收到坐标与振幅中，使新相位精确成为二次型。

\Needspace{6\baselineskip}
\begin{lemma}[Morse 局部归约]\label{b2:lem:21-morse-normal-form}
设整数 \(d\geq1\)，\(U\subset\mathbb R^d\) 开，\(\phi\in C^\infty(U;\mathbb R)\)，\(x_0\in U\) 满足 \(\nabla\phi(x_0)=0\) 且 \(H_0=D^2\phi(x_0)\) 可逆。则存在 \(x_0\) 的邻域 \(U_0\)、原点的邻域 \(V_0\)，以及光滑微分同胚 \(u=\kappa(x):U_0\longrightarrow V_0\)，使 \(\kappa(x_0)=0\) 且
\[
 \phi(x)=\phi(x_0)+\frac12u^{\mathsf T}Su,
 \quad
 S=\operatorname{diag}(\varepsilon_1,\ldots,\varepsilon_d),
 \quad \varepsilon_j\in\{1,-1\}.
\]
其中 \(\sum_j\varepsilon_j=\operatorname{sgn}H_0\)。记逆映射为 \(x=X(u)\)，则
\[
 |\det DX(0)|=|\det H_0|^{-1/2}.
\]
\end{lemma}
\begin{proof}
先平移 \(x_0\) 到原点，再作一个固定正交换元，将 \(H_0\) 化为
\(\operatorname{diag}(\mu_1,\ldots,\mu_d)\)，其中每个 \(\mu_j\neq0\)。这两次换元的 Jacobi 因子绝对值都为一。暂用 \(x\) 表示变换后的坐标。在足够小的球内，线段 \(tx\)、\(0\leq t\leq1\)，都留在相位定义域中。带积分余项的 Taylor 公式给出精确恒等式
\[
 \phi(x)-\phi(0)=\frac12x^{\mathsf T}G(x)x,
 \quad
 G(x)=2\int_0^1(1-t)D^2\phi(tx)\,\mathrm dt.
\]
矩阵 \(G(x)\) 实对称、光滑，且 \(G(0)=H_0\)。这里把高阶项保留在矩阵函数中，并没有把它们丢进误差。

先看二维的配方。若
\[
 G(x)=
 \begin{pmatrix}p(x)&q(x)\\q(x)&r(x)\end{pmatrix},
\]
则在 \(p\neq0\) 的区域有
\[
 x^{\mathsf T}G(x)x
 =p(x)\left(x_1+\frac{q(x)}{p(x)}x_2\right)^2
  +\left(r(x)-\frac{q(x)^2}{p(x)}\right)x_2^2.
\]
在原点，两个新系数分别为 \(\mu_1,\mu_2\)，因而缩小邻域后都非零且符号不变。把它们绝对值的平方根分别乘入括号，就得到两个新的坐标表达式。系数依赖于原变量 \(x\) 并不破坏这个逐点恒等式；剩下要核验的是这些表达式合起来是否真的可作坐标。

一般维数按同一规则逐步配方。第一步以 \(p_1(x)=G_{11}(x)\) 为主元，提出
\[
 p_1(x)\left(x_1+\sum_{k>1}
                     \frac{G_{1k}(x)}{p_1(x)}x_k\right)^2,
\]
余下变量的对称系数矩阵为
\[
 G^{(2)}_{kl}(x)
 =G_{kl}(x)-\frac{G_{k1}(x)G_{1l}(x)}{p_1(x)}
 \quad(k,l>1).
\]
由于 \(G(0)\) 已经对角化，\(p_1(0)=\mu_1\)，且
\(G^{(2)}(0)=\operatorname{diag}(\mu_2,\ldots,\mu_d)\)。下一步以 \(G^{(2)}_{22}\) 为主元，重复同一配方。归纳可知，第 \(j\) 步主元 \(p_j\) 在原点等于 \(\mu_j\)，所以每一步都可以在一个更小邻域内继续。只有有限步，最后取这些邻域的公共部分，所有主元均不为零、符号固定，且全部商与平方根都光滑。

因此得到一个光滑可逆的上三角矩阵函数 \(B(x)\)，使
\[
 G(x)=B(x)^{\mathsf T}S B(x),
 \quad
 S=\operatorname{diag}\bigl(\operatorname{sgn}\mu_1,\ldots,
                            \operatorname{sgn}\mu_d\bigr).
\]
这里 \(B\) 的对角元为 \(\sqrt{|p_j|}\)，其余系数来自上述配方。令
\(u=\kappa(x)=B(x)x\)，则已经有精确的相位恒等式。还需注意：矩阵 \(B(x)\) 逐点可逆，本身不等于映射 \(x\mapsto B(x)x\) 已知可逆。真正的依据是在原点
\[
 D\kappa(0)=B(0),
\]
因为对 \(B\) 求导后乘上的 \(x\) 在原点为零。矩阵 \(B(0)\) 可逆，逆函数定理于是给出光滑局部逆映射。

在驻点处对相位恒等式求两次导数，得到
\[
 H_0=D\kappa(0)^{\mathsf T}S D\kappa(0).
\]
取行列式绝对值，并用 \(|\det S|=1\)，便有
\[
 |\det D\kappa(0)|^2=|\det H_0|,
 \quad
 |\det DX(0)|=|\det H_0|^{-1/2}.
\]
恢复原来的正交坐标后，这一绝对值公式不变。符号差由原点的对角化已经确定，故等于 \(\operatorname{sgn}H_0\)。
\end{proof}

预先对角化的作用是保证固定消元顺序中的主元不消失。仅有矩阵可逆不够：\(\left(\begin{smallmatrix}0&1\\1&0\end{smallmatrix}\right)\) 可逆，却不能直接把左上角作为第一个主元。本证明只在一个固定矩阵上选择正交坐标，随后用连续性工作；没有要求在一个矩阵族中光滑地选择所有特征向量。

\ProseParagraphBreak

例如
\[
 \phi(x_1,x_2)=\frac12\left[x_1^2+2x_1x_2+(2+x_1)x_2^2\right]
 =\frac12\left[(x_1+x_2)^2+(1+x_1)x_2^2\right].
\]
在原点附近可取 \(u_1=x_1+x_2\)、\(u_2=\sqrt{1+x_1}\,x_2\)。其在原点的导数行列式为一，故确为局部坐标，并将含三次项的相位精确化为 \((u_1^2+u_2^2)/2\)。原 Hesse 矩阵为
\(\left(\begin{smallmatrix}1&1\\1&2\end{smallmatrix}\right)\)，行列式也为一，与 Jacobi 公式相符。

\ProseParagraphBreak

换元以后振幅成为 \(a(X(u))|\det DX(u)|\)，不能只替换函数自变量而忽略体积因子。一个 Hesse 特征值变小时，对应方向的驻点尺度变宽，\(|\det H_0|^{-1/2}\) 也可能增大；这是非退化下界必须进入估计常数的原因。

\subsection{驻相引理}
\label{b2:subsec:210203}

二次振荡核没有绝对可积性，不能直接把实 Gaussian 积分中的正参数换成纯虚数。先在实部为正的参数域计算，再取边界极限，便能同时确定常数与复平方根的分支。

\Needspace{6\baselineskip}
\begin{proposition}[复 Gaussian 的 Fourier 变换]\label{b2:prop:21-complex-gaussian}
设 \(a\in\mathbb C\)、\(\operatorname{Re}a>0\)、\(\xi\in\mathbb R\)。则下列积分绝对收敛，且
\[
 F(a,\xi)=\int_{\mathbb R}e^{-\pi a x^2}e^{-2\pi\mathrm ix\xi}\,\mathrm dx
 =a^{-1/2}e^{-\pi\xi^2/a}.
\]
这里 \(a^{-1/2}=|a|^{-1/2}e^{-\mathrm i\arg(a)/2}\)，
\(-\pi/2<\arg(a)<\pi/2\)；在正实轴上取正平方根。
\end{proposition}
\begin{proof}
固定 \(\varepsilon>0\)，写 \(a=\varepsilon+\mathrm it\)。任意有限 \(t,\xi\) 区间上的参数导数由多项式乘 \(e^{-\pi\varepsilon x^2}\) 控制，故可以在积分号下求导；下面的实轴分部积分也没有无穷远边界项。积分
\(\partial_x(e^{-\pi a x^2-2\pi\mathrm ix\xi})\) 得到
\(\int x e^{-\pi a x^2-2\pi\mathrm ix\xi}\,\mathrm dx=-\mathrm i\xi F/a\)，所以
\(\partial_\xi F=-(2\pi\xi/a)F\)。解这个实变量一阶方程可得
\(F(a,\xi)=F(a,0)e^{-\pi\xi^2/a}\)。

再积分 \(\partial_x(xe^{-\pi a x^2})\)，得到
\(\int x^2e^{-\pi a x^2}\,\mathrm dx=F(a,0)/(2\pi a)\)。于是
\(\partial_tF(\varepsilon+\mathrm it,0)=-\mathrm iF/(2a)\)。实 Gaussian 公式给初值 \(F(\varepsilon,0)=\varepsilon^{-1/2}\)，上述连续分支的 \((\varepsilon+\mathrm it)^{-1/2}\) 满足同一方程和初值，故两者相等。整个计算只沿实轴作分部积分，复平方根由初值连续确定。
\end{proof}
由于 \(\operatorname{Re}(1/a)>0\)，此公式还有
\( |F(a,\xi)|\leq |a|^{-1/2}\)。在纯二次振荡的正则化中取
\(a=\varepsilon/\pi-\mathrm i\lambda\varepsilon_j\)、\(\varepsilon_j=\pm1\)，则模上界为 \(\lambda^{-1/2}\)，与正则化参数无关。高维对角二次型逐坐标相乘即可；乘上可积的 \(\widehat b\) 后，这正是下面支配收敛所需的共同控制。

先把二次型与光滑振幅的配对写成一个可绝对积分的 Fourier 表达式。这不仅给出首项，也给出每个更高阶系数及可检查的余项。

\Needspace{6\baselineskip}
\begin{proposition}[二次驻相展开]\label{b2:prop:21-quadratic-expansion}
设整数 \(d\geq1\)，\(S=\operatorname{diag}(\varepsilon_1,\ldots,\varepsilon_d)\)，其中 \(\varepsilon_j\in\{1,-1\}\)，且 \(b\in C_c^\infty(\mathbb R^d;\mathbb C)\)。采用负号 \(2\pi\) Fourier 归一化，并令
\[
 D_S=\sum_{k=1}^d\varepsilon_k\partial_{u_k}^2.
\]
对 \(\lambda\geq1\)，有精确恒等式
\begin{equation}\label{b2:eq:21-quadratic-fourier}
 \int_{\mathbb R^d}e^{\pi\mathrm i\lambda u^{\mathsf T}Su}
                 \cdot b(u)\,\mathrm du
 =\lambda^{-d/2}e^{\mathrm i\pi\operatorname{sgn}S/4}
   \int_{\mathbb R^d}e^{-\pi\mathrm i\eta^{\mathsf T}S\eta/\lambda}
                 \cdot\widehat b(\eta)\,\mathrm d\eta.
\end{equation}
对任意整数 \(N\geq1\)，因而有
\[
 \begin{aligned}
 \int e^{\pi\mathrm i\lambda u^{\mathsf T}Su}\cdot b(u)\,\mathrm du
 &=\lambda^{-d/2}e^{\mathrm i\pi\operatorname{sgn}S/4}
   \left[
   \sum_{j=0}^{N-1}\frac1{j!}
     \left(\frac{\mathrm i}{4\pi\lambda}\right)^j
     D_S^jb(0)+R_N(\lambda)
   \right],\\
 |R_N(\lambda)|
 &\leq\frac{\pi^N}{N!}\lambda^{-N}
      \int_{\mathbb R^d}|\eta|^{2N}\cdot
                        |\widehat b(\eta)|\,\mathrm d\eta.
 \end{aligned}
\]
若 \(\operatorname{supp}b\subset K\)，其中 \(K\) 为固定紧集，则对任意整数 \(M>2N+d\)，右端积分至多为
\[
 C_{K,M,N}\sum_{|\alpha|\leq M}
                    \|\partial^\alpha b\|_\infty.
\]
\end{proposition}
\begin{proof}
先在左侧乘 \(e^{-\epsilon|u|^2}\)，其中 \(\epsilon>0\)。由 Schwartz 函数的 Fourier 反演，
\[
 b(u)=\int_{\mathbb R^d}e^{2\pi\mathrm i u\cdot\eta}
                  \cdot\widehat b(\eta)\,\mathrm d\eta.
\]
阻尼 Gaussian 在 \(u\) 上可积，\(\widehat b\) 在 \(\eta\) 上可积，因此可用 Fubini 交换两个积分。内层积分分成 \(d\) 个一维积分，在\cref{b2:prop:21-complex-gaussian}中分别取
\[
 a_k=\epsilon/\pi-\mathrm i\lambda\varepsilon_k,
 \quad \xi=-\eta_k.
\]
由于 \(\operatorname{Re}(1/a_k)>0\)，每一因子的模不超过
\(|a_k|^{-1/2}\leq\lambda^{-1/2}\)。所以交换后的被积函数统一由
\(\lambda^{-d/2}|\widehat b(\eta)|\) 控制。让 \(\epsilon\downarrow0\) 时，
\[
 a_k^{-1/2}\longrightarrow
        \lambda^{-1/2}e^{\mathrm i\pi\varepsilon_k/4},
 \quad
 e^{-\pi\eta_k^2/a_k}\longrightarrow
        e^{-\pi\mathrm i\varepsilon_k\eta_k^2/\lambda}.
\]
左侧则由可积函数 \(|b|\) 控制。两侧分别使用支配收敛，便得到\cref{b2:eq:21-quadratic-fourier}。Gaussian 阻尼保证交换积分合法，可积控制则使去掉阻尼后的极限仍保留该恒等式。

对实数 \(t\)，Taylor 积分余项给出
\[
 \left|e^{\mathrm it}
       -\sum_{j=0}^{N-1}\frac{(\mathrm it)^j}{j!}\right|
 \leq\frac{|t|^N}{N!}.
\]
在\cref{b2:eq:21-quadratic-fourier}中取
\(t=-\pi\eta^{\mathsf T}S\eta/\lambda\)，并用
\(|\eta^{\mathsf T}S\eta|\leq|\eta|^2\)，得到所列余项界。各个多项式项可以积分，因为 \(\widehat b\) 是 Schwartz 函数。Fourier 求导公式给出
\[
 \int_{\mathbb R^d}(\eta^{\mathsf T}S\eta)^j
                   \cdot\widehat b(\eta)\,\mathrm d\eta
 =(-4\pi^2)^{-j}D_S^jb(0).
\]
将这一因子与 \((-\pi\mathrm i/\lambda)^j\) 相乘，正好得到
\((\mathrm i/(4\pi\lambda))^j\)，从而确定次项的符号和常数。

最后估计所需的 Fourier 矩。对 \(|\eta|\geq1\)，选择
\(|\eta_k|\geq|\eta|/\sqrt d\)，沿第 \(k\) 个坐标分部积分 \(M\) 次，得到
\[
 |\widehat b(\eta)|
 \leq C_{d,M}|\eta|^{-M}\cdot\|\partial_k^Mb\|_1.
\]
对 \(|\eta|\leq1\) 用 \(|\widehat b(\eta)|\leq\|b\|_1\)。固定紧支撑使这些 \(L^1\) 范数由 \(L^\infty\) 半范数控制，故
\[
 |\widehat b(\eta)|
 \leq C_{K,M}(1+|\eta|)^{-M}
       \sum_{|\alpha|\leq M}\|\partial^\alpha b\|_\infty.
\]
条件 \(M>2N+d\) 保证乘上 \(|\eta|^{2N}\) 后仍可积，证明完成。
\end{proof}

下面采用\term[zhuxiang-gongshi]{驻相公式}[stationary phase formula]的下述版本。指数固定为 \(e^{2\pi\mathrm i\lambda\phi}\)：相位中二次项的系数是 \(1/2\)，因此归约后的指数是 \(e^{\pi\mathrm i\lambda u^{\mathsf T}Su}\)。

\Needspace{6\baselineskip}
\begin{theorem}[驻相公式]\label{b2:thm:21-stationary-phase}
设整数 \(d\geq1\)，\(U\subset\mathbb R^d\) 开，\(a\in C_c^\infty(U;\mathbb C)\)，\(\phi\in C^\infty(U;\mathbb R)\)。假定 \(\operatorname{supp}a\) 中只有一个驻点 \(x_0\)，且 \(H=D^2\phi(x_0)\) 可逆。则当 \(\lambda\to+\infty\) 时，
\[
 \int e^{2\pi\mathrm i\lambda\phi(x)}a(x)\,\mathrm dx
 =e^{2\pi\mathrm i\lambda\phi(x_0)}
   \lambda^{-d/2}e^{\mathrm i\pi\operatorname{sgn}H/4}
   \frac{a(x_0)}{|\det H|^{1/2}}
   +O(\lambda^{-d/2-1}).
\]
\end{theorem}
\begin{proof}
由\cref{b2:lem:21-morse-normal-form}，在 \(x_0\) 附近选择 Morse 坐标 \(u=\kappa(x)\)，逆映射记为 \(X\)。取一个支撑于该坐标邻域的光滑截断 \(\chi\)，在 \(x_0\) 的更小邻域内等于一，分解
\(a=\chi a+(1-\chi)a\)。第二部分的紧支撑上没有驻点，所以梯度有正下界；\cref{b2:prop:21-nonstationary-multi}使其积分为 \(O(\lambda^{-M})\)，其中整数 \(M\) 可以任意选取。

第一部分换元后成为
\[
 e^{2\pi\mathrm i\lambda\phi(x_0)}
 \int e^{\pi\mathrm i\lambda u^{\mathsf T}Su}
                   \cdot b(u)\,\mathrm du,
 \quad
 b(u)=(\chi a)(X(u))\cdot|\det DX(u)|.
\]
由于 \(\chi a\) 的支撑紧包含于坐标邻域，\(b\) 向坐标像外延零后属于 \(C_c^\infty(\mathbb R^d)\)，没有人为造成的边界跳跃。Morse 引理给出
\[
 b(0)=a(x_0)|\det H|^{-1/2},
 \quad
 \operatorname{sgn}S=\operatorname{sgn}H.
\]
对\cref{b2:prop:21-quadratic-expansion}取 \(N=1\)，再在非驻相部分取 \(M>d/2+1\)，便得到所述主项和余项。
\end{proof}
二次模型的误差已由有限个振幅导数控制。Morse 换元把 \(b\) 写成原振幅与 Jacobi 因子的乘积，有限次链式法则又将这些导数化成相位、振幅及逆映射的有限阶导数。因此参数族的统一性需要控制支撑、Hesse 最小奇异值、所用导数，以及远离驻点部分的梯度下界；单独记录每个参数的行列式非零不够。

该公式是固定振幅、固定相位下的渐近式。若振幅支撑随 \(\lambda\) 缩小或导数增长，余项常数可能增长；此时应先把自然尺度规范化。

% 扩写定位：b2:subsec:210203
边界还会改变驻相的首项。考虑一个临界点恰好位于积分端点的模型
\[
 J(\lambda)=\int_0^\infty e^{\pi\mathrm i\lambda t^2}a(t)\,\mathrm dt,
 \quad a\in C_c^\infty([0,\infty)).
\]
此处光滑是指可延拓到零点邻域的光滑函数。若 \(a(0)\ne0\)，则
\[
 J(\lambda)=\frac12e^{\pi\mathrm i/4}a(0)\lambda^{-1/2}+O(\lambda^{-1}).
\]
取在零点附近为一的偶截断 \(\chi\in C_c^\infty(\mathbb R)\)，在半直线上分解
\[
 a(t)=a(0)\chi(t)+t b(t).
\]
分子在零点消失，商 \(b(t)=[a(t)-a(0)\chi(t)]/t\) 因而可光滑延拓到零点，并且仍紧支撑。第一项的半直线积分由偶性恰为全直线截断二次积分的一半；\cref{b2:prop:21-quadratic-expansion}给出它的主项
\(\tfrac12e^{\pi\mathrm i/4}a(0)\lambda^{-1/2}\)。
第二项令 \(v=t^2\)，成为
\[
 \frac12\int_0^\infty e^{\pi\mathrm i\lambda v}
                         \cdot b(\sqrt v)\,\mathrm dv.
\]
新振幅在零点连续，且对 \(v>0\)，其导数为
\(b'(\sqrt v)/(2\sqrt v)\)，绝对积分恰为
\(\int_0^\infty|b'(t)|\,\mathrm dt\)。因此它绝对连续，虽有可能的 \(v^{-1/2}\) 导数奇性，仍可作一阶分部积分，得到界
\[
 \frac1{2\pi\lambda}
       \left[|b(0)|+\int_0^\infty|b'(t)|\,\mathrm dt\right].
\]
这证明了余项为 \(O(\lambda^{-1})\)，也确定了首项中的一半来自端点几何，而不是 Fourier 归一化的改变。

\ProseParagraphBreak

若临界点离开端点而进入内部，其首项重新是完整的二次驻相系数；若移到区间外，则固定参数下没有临界点而只剩边界项。随着临界点在距离 \(\lambda^{-1/2}\) 的区域内移动，应使用保留端点的 Fresnel 积分，不能把内部公式和非驻相公式硬拼成一个统一展开。这个模型同时说明，振幅紧支于开集和振幅在闭区间端点非零是不同假设。

\subsection{渐近展开}
\label{b2:subsec:210204}

\Cref{b2:prop:21-quadratic-expansion}已经给出二次模型的任意阶展开。一般光滑相位先经 Morse 换元，每个系数便由新振幅 \(b\) 在原点的有限阶导数确定：第 \(j\) 项用到 \(D_S^jb(0)\)，因而需要新振幅的至多 \(2j\) 阶导数。若相位原本就是 \(x^{\mathsf T}Hx/2\)，其中 \(H\) 实对称可逆，则只需固定线性换元。这时记
\[
 D_H=\sum_{k,l=1}^d(H^{-1})_{kl}\partial_{x_k}\partial_{x_l},
\]
可将前两项直接写成
\[
 \begin{aligned}
 \int e^{\pi\mathrm i\lambda x^{\mathsf T}Hx}
                  \cdot a(x)\,\mathrm dx
 ={}&\frac{\lambda^{-d/2}e^{\pi\mathrm i\operatorname{sgn}H/4}}
             {|\det H|^{1/2}}\\
 &{}\times
 \left[a(0)+\frac{\mathrm i}{4\pi\lambda}D_Ha(0)
               +O(\lambda^{-2})\right],
 \quad a\in C_c^\infty(\mathbb R^d).
 \end{aligned}
\]
具体地，若 \(x=Tu\) 且 \(T^{\mathsf T}HT=S\)，则
\(H^{-1}=TST^{\mathsf T}\)。链式法则把 \(D_S\) 化为上面的 \(D_H\)，线性换元的 Jacobi 因子为常数 \(|\det H|^{-1/2}\)。这一步不能直接推广为非二次相位的次项公式，因为非线性换元的 Jacobi 因子也会变化。

\ProseParagraphBreak
领先系数为零并不意味着整个积分为零。例如一维二次相位配以偶振幅 \(a(t)=t^2\chi(t)\)，其中 \(\chi\in C_c^\infty(\mathbb R)\) 偶且在零点附近为一。此时 \(a(0)=0\)、\(a''(0)=2\)，所以
\[
 \int e^{\pi\mathrm i\lambda t^2}\cdot t^2\chi(t)\,\mathrm dt
 =\frac{\mathrm i e^{\pi\mathrm i/4}}{2\pi}
        \lambda^{-3/2}+O(\lambda^{-5/2}).
\]
若改用奇振幅 \(t\chi(t)\)，积分则由对称性恒为零。前者是局部消失阶带来的增益，后者是全局对称性抵消；两种机制应分清。

\ProseParagraphBreak

实际应用常把 Fourier 方向、曲面上的位置或另一个积分变量当作参数。此时逐个参数都能使用驻相还不够，还须知道邻域、余项及其参数导数可以统一选择。下面给出一个固定紧支撑的版本。

\Needspace{6\baselineskip}
\begin{proposition}[含参数的驻相展开]\label{b2:prop:21-parametric-stationary}
设整数 \(d,m\geq1\)，\(U\subset\mathbb R^d\)、\(V\subset\mathbb R^m\) 开，\(K\Subset U\) 为紧集。设
\(\phi\in C^\infty(U\times V;\mathbb R)\)、
\(a\in C^\infty(U\times V;\mathbb C)\)，且每个 \(v\in V\) 都满足
\(\operatorname{supp}_x a(\,\cdot\,,v)\subset K\)。假定对每个 \(v\in V\)，\(\phi(\,\cdot\,,v)\) 在 \(K\) 中恰有一个驻点 \(x_0(v)\)，并且
\[
 H(v)=D_x^2\phi(x_0(v),v)
\]
可逆。令
\[
 c(v)=\phi(x_0(v),v),
 \quad
 I(\lambda,v)=\int_U e^{2\pi\mathrm i\lambda\phi(x,v)}
                              \cdot a(x,v)\,\mathrm dx.
\]
则 \(x_0,c\) 光滑，存在光滑系数 \(c_j(v)\)，使对任意整数 \(N\geq1\)、
紧集 \(P\Subset V\) 及参数多重指标 \(\beta\)，
\[
 \begin{aligned}
 e^{-2\pi\mathrm i\lambda c(v)}\lambda^{d/2}I(\lambda,v)
 &=\sum_{j=0}^{N-1}\lambda^{-j}c_j(v)+R_N(\lambda,v),\\
 \sup_{v\in P}|\partial_v^\beta R_N(\lambda,v)|
 &\leq C_{P,N,\beta}\lambda^{-N}
 \quad(\lambda\geq1).
 \end{aligned}
\]
首项为
\[
 c_0(v)=e^{\pi\mathrm i\operatorname{sgn}H(v)/4}
           \frac{a(x_0(v),v)}{|\det H(v)|^{1/2}}.
\]
\end{proposition}
\begin{proof}
固定 \(\lambda\) 时，公共紧支撑保证可以在积分号下对参数任意次求导。先核验驻点随参数连续。若 \(v_n\to v\)，则 \(x_0(v_n)\in K\) 的任意收敛子列极限仍在 \(K\)，梯度的连续性使它成为参数 \(v\) 的驻点。唯一性迫使这个极限为 \(x_0(v)\)。紧性于是给出整个驻点列的收敛。在
\(\nabla_x\phi(x,v)=0\) 上应用含参数的隐函数定理，由 \(H(v)\) 可逆得到 \(x_0(v)\) 光滑，\(c(v)\) 也随之光滑。

固定 \(v_*\in V\)，先把 \(x_0(v)\) 平移到原点。选择一个固定正交矩阵将 \(H(v_*)\) 对角化，而不随 \(v\) 选择特征向量。参数接近 \(v_*\) 且空间变量足够小时，\cref{b2:lem:21-morse-normal-form}证明中的全部消元主元仍非零且符号固定。因此同一配方给出光滑依赖 \((x,v)\) 的 Morse 映射。含参数的逆函数定理使逆映射 \(X(u,v)\) 同样光滑，二次矩阵 \(S\) 在这一参数片上固定。

对紧参数集 \(P\) 取有限个这样的参数片，并在每片中选随 \(x_0(v)\) 平移的截断。必要时缩小各片，使截断后的新振幅支撑都包含于一个固定 \(u\) 紧集，而且每个预先指定阶的 \(u,v\) 导数一致有界。这来自光滑映射在稍小的紧集上的有界性。在临界点处，Hesse 最小奇异值的正下界保证逆映射的一阶导数受控；反复微分逆映射恒等式，便把更高阶导数界归结为相位的相应有限阶导数。

去掉临界值因子后，局部积分正是一个固定二次相位配光滑参数振幅 \(b(u,v)\)。对\cref{b2:prop:21-quadratic-expansion}及其余项估计应用 \(\partial_v^\beta\)，只需把 \(b\) 换成 \(\partial_v^\beta b\)。于是每片上的系数光滑，余项连同该阶参数导数均为 \(O(\lambda^{-N})\)。在重叠参数片上，系数相同：将两种展开相减，先令 \(\lambda\to\infty\) 比较常数项，再依次乘以 \(\lambda\) 比较下一项，便逐项得到唯一性。因此这些局部系数可拼成 \(V\) 上的光滑函数。

还须控制截断之外的部分。它在 \(K\times P\) 内远离驻点图像，且没有其他驻点，故全梯度 \(|\nabla_x\phi|\) 有统一正下界。对归一化后的非驻相积分求 \(\beta\) 阶参数导数，最多产生 \(|\beta|\) 个 \(\lambda\) 因子，外面另有 \(\lambda^{d/2}\)。预先将\cref{b2:prop:21-nonstationary-multi}中的积分分部次数取为大于 \(N+|\beta|+d/2\) 的整数，这一部分仍为 \(O(\lambda^{-N})\)。有限个参数片上的常数取最大值，即得结论。
\end{proof}

命题中的紧性把非退化条件转为统一下界。若参数范围不紧，或相位族不在公共邻域光滑，必须另行假定固定支撑、Hesse 的统一非退化下界、所需有限阶导数界，以及远离驻点区域的梯度下界。特别要先去掉迅速变化的临界值因子：直接微分 \(e^{2\pi\mathrm i\lambda c(v)}\) 会产生 \(\lambda\)，不能把未归一化积分的参数导数误当作与函数值具有同一衰减阶。

% 扩写定位：b2:subsec:210204
驻相公式也适用于一个临界子流形，前提是垂直于临界片的二次型非退化。在坐标已经把临界片拉直为 \(z=0\) 时，可以把切向变量 \(y\) 当作参数，再在法向变量 \(z\) 上应用刚才的统一结论。若这些坐标来自先前的一次换元，体积的 Jacobi 因子须已经吸收到下面的振幅 \(a\) 中。

\Needspace{6\baselineskip}
\begin{proposition}[非退化法向临界片]\label{b2:prop:21-clean-stationary}
设整数 \(d\geq1\)、\(0\leq r\leq d\)，\(Y\subset\mathbb R^{d-r}\) 是非空连通开集，\(Z\subset\mathbb R^r\) 是包含原点的开集。设
\(\phi\in C^\infty(Y\times Z;\mathbb R)\)、
\(a\in C_c^\infty(Y\times Z;\mathbb C)\)，且对所有 \(y\in Y\) 都有
\(\mathrm d\phi(y,0)=0\)。假定 \(\operatorname{supp}a\) 中除 \(z=0\) 外没有其他全梯度驻点。当 \(r\geq1\) 时，进一步假定法向 Hesse 矩阵
\(H(y)=\partial_z^2\phi(y,0)\) 对所有 \(y\in Y\) 可逆，其最小奇异值在振幅的 \(y\) 投影附近有统一正下界。则 \(\phi(y,0)=c\) 为常数，且当 \(\lambda\to+\infty\) 时，
\[
 \int e^{2\pi\mathrm i\lambda\phi(y,z)}a(y,z)\,\mathrm dy\,\mathrm dz
 =e^{2\pi\mathrm i\lambda c}\lambda^{-r/2}
 \int e^{\pi\mathrm i\operatorname{sgn}H(y)/4}
 \frac{a(y,0)}{|\det H(y)|^{1/2}}\,\mathrm dy
 +O(\lambda^{-r/2-1}).
\]
约定 \(\mathbb R^0=\{0\}\)，零维积分取函数在该点的值。\(r=0\) 时取空矩阵行列式为一、符号差为零，公式成为没有余项的恒等式；\(r=d\) 时没有 \(y\) 变量，公式即孤立驻点情形。
\end{proposition}
\begin{proof}
临界条件使函数 \(y\mapsto\phi(y,0)\) 的各一阶导数为零，故它局部常值；连通性使其在 \(Y\) 上为同一个常数。若 \(r=0\)，全部变量都是切向变量，相位处处为 \(c\)，直接提出指数因子即得恒等式。若 \(r=d\) 且原点在振幅支撑中，命题就是前面的驻相公式；若原点不在支撑中，则主项为零，\cref{b2:prop:21-nonstationary-multi}给出任意阶衰减。以下设 \(1\leq r<d\)。

振幅的 \(y\) 投影为紧集。在其中每一点附近，将 \(y\) 当作参数，\(\partial_z\phi(y,0)=0\) 与 \(H(y)\) 可逆保证 \(z=0\) 在一个共同的小法向邻域内是唯一的法向驻点。取有限个这样的参数片与光滑分割，并在各片内加一个支撑于该法向邻域、在零点附近为一的截断。\Cref{b2:prop:21-parametric-stationary}于是适用于每片的 \(z\) 积分，给出
\[
 e^{2\pi\mathrm i\lambda c}\lambda^{-r/2}
 e^{\pi\mathrm i\operatorname{sgn}H(y)/4}
 \frac{a_{\mathrm{loc}}(y,0)}{|\det H(y)|^{1/2}}
 +O(\lambda^{-r/2-1}),
\]
其中余项对该片的 \(y\) 一致。它们的 \(y\) 支撑均在固定紧集中，故可对 \(y\) 积分，余项保持同阶；在 \(z=0\) 上各截断之和为一，主项相加就是题中积分。

剩余振幅紧支撑且远离 \(z=0\)。题设排除了此处的全梯度驻点，紧性给出 \(|\nabla_{y,z}\phi|\) 的正下界。用\cref{b2:prop:21-nonstationary-multi}在全部 \(d\) 个变量上分部积分，可取得任意阶衰减。这里不能仅凭全梯度非零，就对每个固定 \(y\) 在 \(z\) 中分部积分：远处可能有 \(\partial_z\phi=0\) 而 \(\partial_y\phi\neq0\)。先局部处理法向驻点，再对余部使用全梯度，才覆盖了题设允许的情况。
\end{proof}
衰减幂次由真正振荡的法向维数 \(r\) 决定，而不是积分空间的总维数。例如 \(\phi(y,z)=|z|^2/2\) 与乘积振幅 \(a=b(y)c(z)\)，对 \(y\) 的积分没有任何 \(\lambda\) 衰减；只有 \(z\) 变量贡献 \(\lambda^{-r/2}\)。圆柱曲面的 Fourier 变换也体现这种秩信息：弯曲方向提供驻相增益，平直方向是否衰减取决于振幅和 Fourier 方向，不能从总维数套用球面的公式。

\ProseParagraphBreak

这一局部公式不是任意退化临界点的公式。若 Hesse 矩阵在法向也退化，Morse 坐标不再存在；若临界集有交叉或尖点，也不能选成单一光滑片。比如 \(t^4\) 在零点的衰减为 \(\lambda^{-1/4}\)，已与任何一维非退化二次型的 \(\lambda^{-1/2}\) 不同。所需的是有限型分析，而不是把行列式的零值放进分母。

\ProseParagraphBreak

非二次相位的下一项并不只取决于振幅的二阶导数。这个差别可以用一维局部计算明确看到。为固定常数，假设
\(\phi(0)=\phi'(0)=0\)、\(\phi''(0)=1\)，记
\(c_3=\phi'''(0)\)、\(c_4=\phi^{(4)}(0)\)，于是
\[
 \phi(t)=\frac{t^2}2+\frac{c_3}6t^3+\frac{c_4}{24}t^4+O(t^5).
\]
商 \(2\phi(t)/t^2\) 在零点有光滑延拓且值为一，所以在足够小的邻域取正平方根，定义
\[
 y=t\sqrt{\frac{2\phi(t)}{t^2}}
   =t+\frac{c_3}{6}t^2
       +\left(\frac{c_4}{24}-\frac{c_3^2}{72}\right)t^3+O(t^4).
\]
其导数在零点等于一，故有光滑逆映射。把待定的三次 Taylor 多项式代回上式、逐项比较，得到
\[
 t=y-\frac{c_3}6y^2+
           \left(\frac{5c_3^2}{72}-\frac{c_4}{24}\right)y^3+O(y^4).
\]
逆映射在此保持方向，因此新振幅是
\(b(y)=a(t(y))\cdot t'(y)\)。若记逆映射的二次、三次系数分别为 \(\alpha,\beta\)，则 \(b\) 的二次系数是
\(a''(0)/2+3\alpha a'(0)+3\beta a(0)\)。代入这两个系数，得到
\[
 b''(0)=a''(0)-c_3a'(0)
 +\left(\frac{5c_3^2}{12}-\frac{c_4}4\right)a(0).
\]
二次模型的展开于是给出
\[
 I_\lambda=e^{\pi\mathrm i/4}\lambda^{-1/2}
 \left[a(0)+\frac{\mathrm i}{4\pi\lambda}b''(0)
 +O(\lambda^{-2})\right],
\]
这里振幅光滑紧支于该局部片，或已用截断把其余非驻相部分分开。若未减去原来的临界值，整个右侧还须乘以 \(e^{2\pi\mathrm i\lambda\phi(0)}\)。

\ProseParagraphBreak

即使原振幅在驻点附近为常数，三次和四次相位也会改变次项，因为坐标变换的体积因子不是常数。若只把二次公式中的 \(a''(0)\) 原样用于任意相位，就会遗漏这些项。计算首项时只需变换在驻点的一阶微分；计算更高项则需要相位与变换的更高阶 Taylor 数据。这解释了为什么完整渐近展开的导数假设随所需阶数增加。

\ProseParagraphBreak

有限个驻点的贡献需要先相加，才能讨论积分的整体大小。假定支撑中只有非退化驻点 \(x_1,\ldots,x_\ell\)，选各点附近互不相交的截断，再把其余部分留给非驻相法，便有
\[
 I_\lambda=\lambda^{-d/2}\sum_{\nu=1}^{\ell}
 e^{2\pi\mathrm i\lambda\phi(x_\nu)}
 e^{\pi\mathrm i\operatorname{sgn}H_\nu/4}
 \frac{a(x_\nu)}{|\det H_\nu|^{1/2}}
 +O(\lambda^{-d/2-1}).
\]
这里每个局部首项都可能非零，但它们的和仍可能抵消。取一维相位 \(\phi(t)=t^3/3-t\)，振幅只支在 \(\pm1\) 的两个小邻域，并在两个驻点等于一。两点的相位值为 \(\mp2/3\)，二阶导数为 \(\pm2\)，故首项合成为
\[
 \sqrt2\,\lambda^{-1/2}
 \cos\left(\frac{4\pi\lambda}{3}-\frac\pi4\right).
\]
沿余弦消失的一列正参数，积分只剩 \(O(\lambda^{-3/2})\)。因此一般驻相定理提供上界与展开；若要求统一的下界，还必须排除不同驻点之间的干涉。反过来，只有一个驻点且该点振幅非零时，首项模有正下界，余项便保证充分大参数下确实具有 \(\lambda^{-d/2}\) 的大小。

\ProseParagraphBreak

首项中的行列式与体积因子也有坐标无关的含义。设 \(x=F(y)\) 是局部光滑换元，\(F(y_0)=x_0\)，且 \(B=DF(y_0)\) 可逆。在驻点处一阶相位导数为零，二阶链式法则中的变换二阶导数项消失，因而新 Hesse 矩阵为 \(B^{\mathsf T}HB\)，新振幅为 \(a(x_0)|\det B|\)。于是
\[
 \frac{a(x_0)|\det B|}
 {|\det(B^{\mathsf T}HB)|^{1/2}}
 =\frac{a(x_0)}{|\det H|^{1/2}}.
\]
实可逆合同也保持正负惯性指数，所以相位因子不变。这说明驻相系数与选择哪一套 Morse 坐标无关。若只改变相位而忘记变换振幅中的 Jacobi 因子，就会破坏这种不变性；前面的一维次项计算正是这一原则在更高 Taylor 阶上的体现。

% !TeX root = ../../Book2.tex
\section{多维振荡积分}
\label{b2:sec:2103}

多维积分中，可利用的振荡方向数未必等于变量个数。相位梯度不为零时，只需沿一个合适的向量场积分分部；非退化驻点附近，每个二次方向提供半个幂次的衰减；若某些方向完全平直，它们只贡献普通积分。本节把这些方法用于参数振幅和曲面 Fourier 变换，并说明统一估计需要怎样局部化。

\subsection{相位函数}
\label{b2:subsec:210301}

\cref{b2:prop:21-nonstationary-multi}已经建立多维非驻相估计。恢复本处所需条件：相位 \(\phi\) 是开集 \(U\subset\mathbb R^d\) 上的实光滑函数，振幅 \(a\in C_c^\infty(U;\mathbb C)\)，并且 \(|\nabla\phi|\geq c>0\) 在振幅支撑的一个邻域成立。相应的向量场及微分算子为
\[
 V(x)=\frac{\nabla\phi(x)}{|\nabla\phi(x)|^2},
 \quad
 L=\frac1{2\pi\mathrm i\lambda}V\cdot\nabla,
 \quad
 Le^{2\pi\mathrm i\lambda\phi}
 =e^{2\pi\mathrm i\lambda\phi}.
\]
在这里转移的是双线性积分中的微分算子，而不是带复共轭的 Hilbert 伴随。由散度公式，其形式转置作用于振幅为
\[
 L^{\mathrm t}a
 =-\frac1{2\pi\mathrm i\lambda}\operatorname{div}(Va)
 =-\frac1{2\pi\mathrm i\lambda}
   \bigl(V\cdot\nabla a+(\operatorname{div}V)a\bigr).
\]
因此即使原振幅局部为常数，相位的变化也会通过 \(\operatorname{div}V\) 进入下一次振幅。具体地，
\[
 \operatorname{div}V
 =\frac{\Delta\phi}{|\nabla\phi|^2}
  -2\,\frac{(D^2\phi\,\nabla\phi)\cdot\nabla\phi}
               {|\nabla\phi|^4}.
\]
梯度下界控制分母；相位的有限阶导数控制分子。反复转移 \(N\) 次后，得到 \(\lambda^{-N}\) 乘以振幅至多 \(N\) 阶导数的有限线性组合，其系数由相位至多 \(N+1\) 阶导数与梯度倒数控制。

\ProseParagraphBreak

这个描述也给出使用参数族时的核对方法。若 \(z\) 遍历某个参数集合，\(a_z\) 的支撑都在同一紧集 \(K\Subset U\) 内，所需导数的 \(L^1\) 范数一致有界，\(\phi_z\) 在同一邻域上满足统一梯度下界和相位导数界，那么任意固定次数的非驻相估计就对 \(z\) 一致。参数集合是否紧只是获得这些界的一种途径；实际使用的是这些定量条件。

\ProseParagraphBreak

紧支撑负责消除积分域边界，不能被“相位在内部没有驻点”替代。对光滑有界区域 \(\Omega\)，一次散度积分分部会出现
\[
 \begin{aligned}
 \int_\Omega e^{2\pi\mathrm i\lambda\phi}\cdot a\,\mathrm dx
 &=
 \frac1{2\pi\mathrm i\lambda}
 \int_{\partial\Omega}
 e^{2\pi\mathrm i\lambda\phi}\cdot a(V\cdot n)\,\mathrm d\sigma\\
 &\hspace{1em}
 -\frac1{2\pi\mathrm i\lambda}
 \int_\Omega e^{2\pi\mathrm i\lambda\phi}\cdot
                    \operatorname{div}(Va)\,\mathrm dx,
 \end{aligned}
\]
其中 \(n\) 是外单位法向。边界积分是否衰减，取决于相位限制在边界后的切向变化。平直立方体上的线性相位甚至可以在某些面上恒定。例如
\[
 \int_{[0,1]^d}e^{2\pi\mathrm i\lambda x_1}\,\mathrm dx
 =\frac{e^{2\pi\mathrm i\lambda}-1}{2\pi\mathrm i\lambda}.
\]
其余 \(d-1\) 个变量只提供体积一，并没有增加衰减。这与光滑紧支撑振幅的快速衰减没有矛盾。

\subsection{振幅}
\label{b2:subsec:210302}

振幅的支撑、导数和在驻点处的消失阶都参与估计。尤其当振幅依赖于 \(\lambda\) 时，不能把它当作固定函数代入渐近公式，再把增长中的导数藏在余项常数里。一个完整可算的模型是
\[
 I_\lambda
 =\int_{\mathbb R^d}e^{\pi\mathrm i\lambda|x|^2}
                    \cdot a(\lambda^\theta x)\,\mathrm dx,
 \quad
 a\in C_c^\infty(\mathbb R^d;\mathbb C),
 \quad \theta\geq0.
\]
这里 \(\lambda\geq1\)，\(\theta\) 固定。换元 \(u=\lambda^\theta x\) 得
\[
 I_\lambda=\lambda^{-d\theta}
 \int_{\mathbb R^d}e^{\pi\mathrm i\tau|u|^2}
                       \cdot a(u)\,\mathrm du,
 \quad
 \tau=\lambda^{1-2\theta}.
\]
振幅的尺度归一化后，真正需要比较的是 \(\tau\) 与一的大小。

\ProseParagraphBreak

当 \(0\leq\theta<1/2\) 时，\(\tau\to\infty\)。在原点附近用二次驻相公式，其余部分用非驻相估计，得到
\[
 I_\lambda
 =e^{\pi\mathrm i d/4}\lambda^{-d/2}
   \left[a(0)+O\bigl(\lambda^{-(1-2\theta)}\bigr)\right].
\]
因子 \(\lambda^{-d\theta}\tau^{-d/2}\) 恰好等于 \(\lambda^{-d/2}\)。此时振幅变化尺度比驻相尺度宽，首项仍只看 \(a(0)\)，但相对于首项的余项是 \(\lambda^{-(1-2\theta)}\)，已不是固定振幅时的 \(\lambda^{-1}\)。若 \(a(0)=0\)，须继续看更高阶项，不能从此式断言首项非零。

\ProseParagraphBreak

当 \(\theta=1/2\) 时，\(\tau=1\)，于是有精确等式
\[
 I_\lambda
 =\lambda^{-d/2}
   \int_{\mathbb R^d}e^{\pi\mathrm i|u|^2}
                       \cdot a(u)\,\mathrm du.
\]
系数依赖整个振幅，通常不能用 \(e^{\pi\mathrm i d/4}a(0)\) 代替。振幅与驻点邻域在同一尺度变化，归一化后的积分已没有一个趋于无穷的参数。

\ProseParagraphBreak

当 \(\theta>1/2\) 时，\(\tau\to0\)。利用
\(|e^{\mathrm it}-1|\leq|t|\) 与 \(a\) 的紧支撑，得到
\[
 I_\lambda
 =\lambda^{-d\theta}
   \left[
   \int_{\mathbb R^d}a(u)\,\mathrm du
   +O\bigl(\lambda^{1-2\theta}\bigr)
   \right].
\]
这里误差由
\(\pi\tau\int|u|^2\cdot|a(u)|\,\mathrm du\)
控制。衰减主要来自支撑体积缩小，而不是相位在该支撑上产生越来越多的波动。这三个范围也说明，同一条固定振幅驻相公式不能在支撑尺度穿过 \(\lambda^{-1/2}\) 时原样保持。

\ProseParagraphBreak

非驻相区域也有类似的导数代价。设 \(\delta\geq0\) 固定。若相位固定、梯度在固定紧集附近有正下界，而 \(a_\lambda\) 支撑于该紧集，并满足
\[
 \|\partial^\gamma a_\lambda\|_1
 \leq A_\gamma\lambda^{\delta|\gamma|},
 \quad \lambda\geq1,
\]
则转移 \(N\) 次导数后至多得到
\[
 |I_\lambda(a_\lambda,\phi)|
 \leq C_N\lambda^{-N(1-\delta)}.
\]
当 \(0\leq\delta<1\)，且每个所需阶数都有上述控制时，仍能通过增加 \(N\) 得到任意固定的负幂。若 \(\delta\geq1\)，这套估计本身不再提供衰减；例如振幅中包含 \(e^{-2\pi\mathrm i\lambda\phi}\) 时，原来的振荡可以被完全消去。

\subsection{曲率}
\label{b2:subsec:210303}

曲面 Fourier 变换把几何曲率转成相位的二阶导数。先在一个图形坐标片中写
\[
 F(y)=(y,h(y)),
 \quad y\in U\subset\mathbb R^{d-1},
 \quad h\in C^\infty(U;\mathbb R).
\]
曲面体积元为
\[
 \mathrm d\sigma(F(y))
 =J(y)\,\mathrm dy,
 \quad
 J(y)=\sqrt{1+|\nabla h(y)|^2}.
\]
这个因子来自切向基向量
\(\partial_jF=e_j+(\partial_jh)e_d\)
的 Gram 矩阵 \(G=I+\nabla h(\nabla h)^{\mathsf T}\)：其行列式为 \(1+|\nabla h|^2\)，体积元是行列式的平方根。忽略这个因子会改变曲面测度，也会改变驻相系数。

\ProseParagraphBreak

取上向单位法向
\(n=(-\nabla h,1)/J\)。
第二基本形式在此切向基下的矩阵为 \(B=D^2h/J\)，因为
\(\partial_{ij}F\cdot n=(\partial_{ij}h)/J\)。
Gaussian 曲率是形算子 \(G^{-1}B\) 的行列式，故
\[
 \mathcal K(y)
 =\frac{\det D^2h(y)}
        {\bigl(1+|\nabla h(y)|^2\bigr)^{(d+1)/2}}.
\]
法向选择可能改变曲率的符号约定，但是否为零及其绝对值不受影响。下面只用其非零性与绝对值。

\Needspace{6\baselineskip}
\begin{proposition}[非零曲率曲面密度的 Fourier 衰减]\label{b2:prop:21-curved-surface-decay}
设 \(d\geq2\) 为整数，\(\Sigma\subset\mathbb R^d\) 是光滑嵌入超曲面，\(a\in C_c^\infty(\Sigma;\mathbb C)\)，且 \(\Sigma\) 在 \(\operatorname{supp}a\) 附近没有边界。假设 \(\Sigma\) 的 Gaussian 曲率在 \(\operatorname{supp}a\) 上处处非零。定义有限复测度 \(\mu=a\,\mathrm d\sigma\)，其中 \(\mathrm d\sigma\) 为诱导曲面测度，并取
\[
 \widehat\mu(\xi)
 =\int_\Sigma e^{-2\pi\mathrm i x\cdot\xi}
                  \cdot a(x)\,\mathrm d\sigma(x).
\]
则存在仅依赖曲面在该紧支撑附近的几何及振幅的常数 \(C\)，使
\[
 |\widehat\mu(\xi)|
 \leq C(1+|\xi|)^{-(d-1)/2},
 \quad \xi\in\mathbb R^d.
\]
\end{proposition}
\begin{proof}
由支撑的紧性，可用有限个经过刚性运动得到的图形坐标片覆盖，并取其上的光滑分割。曲率连续且在支撑上不为零，故可将坐标片略微缩小，使相应紧集附近仍有非零曲率。刚性运动中的平移只给 Fourier 变换乘一个模为一的因子，正交变换只改变方向。因此只需估计每个图形片上的积分
\[
 \int_U e^{2\pi\mathrm i\rho\phi_\omega(y)}
                   \cdot b(y)\,\mathrm dy,
 \quad
 \phi_\omega(y)=-\omega'\cdot y-\omega_dh(y),
\]
其中 \(\rho=|\xi|\)、\(\omega=\xi/|\xi|\in S^{d-1}\)，
\[
 b(y)=\eta(F(y))\cdot a(F(y))\cdot J(y)
\]
是光滑紧支撑振幅，\(\eta\) 为所选分割函数。记 \(K=\operatorname{supp}b\Subset U\)。

先检查任意方向上的驻点。由
\[
 \nabla_y\phi_\omega
 =-\omega'-\omega_d\nabla h,
 \quad
 D_y^2\phi_\omega=-\omega_dD^2h,
\]
若 \(\nabla_y\phi_\omega(y)=0\)，则
\(\omega'=-\omega_d\nabla h(y)\)。
因为 \(|\omega|=1\)，必有
\[
 |\omega_d|=\frac1{J(y)}>0.
\]
这正说明 \(\omega\) 平行于曲面法向。非零曲率又使 \(D^2h(y)\) 可逆，所以驻点处 \(D_y^2\phi_\omega\) 可逆。

为了使常数对方向一致，需要同时在位置和方向上局部化。考虑紧集 \(K\times S^{d-1}\)。对其中每一对 \((y_0,\omega_0)\)，分两种情形。若相位梯度不为零，连续性给出一个乘积邻域，在其中
\(|\nabla_y\phi_\omega|\geq c_0>0\)。
若它为零，先在方向球面上取局部坐标，再由刚才的 Hesse 矩阵可逆性对方程
\(\nabla_y\phi_\omega(y)=0\)
应用含参数的隐函数定理：在足够小的方向邻域内，驻点 \(y_*(\omega)\) 光滑依赖 \(\omega\)，并在一个固定的位置邻域内唯一。进一步缩小邻域，可使其 Hesse 矩阵的最小奇异值有统一正下界，同时保证所有用到的相位导数一致有界。还可在此位置邻域内选一个共同紧集，既包含后面分割振幅的支撑，也包含这些方向下的全部驻点 \(y_*(\omega)\)。

上述乘积邻域覆盖 \(K\times S^{d-1}\)，紧性给出有限子覆盖。可在各片中取位置和方向的光滑截断之积，再除以其总和，构造该紧集附近的有限光滑分割
\(\sum_\nu\chi_\nu(y,\omega)=1\)。
有限性及紧支撑保证分割函数所需的 \(y\) 导数一致有界。于是 \(b\chi_\nu\) 仍支撑于固定紧集，而其导数界不随 \(\omega\) 改变。

对第一类片，\cref{b2:prop:21-nonstationary-multi}给出任意固定阶的 \(O(\rho^{-N})\)，常数对方向一致。对第二类片，\cref{b2:prop:21-parametric-stationary}适用于移动驻点 \(y_*(\omega)\)：固定位置邻域、统一可逆的 Hesse 矩阵、相位与振幅的有限阶导数界均已建立。因此这些片的积分一致为 \(O(\rho^{-(d-1)/2})\)。驻点即使落在某个分割振幅的零区，其首项也只会消失，不影响此上界。取 \(N\geq(d-1)/2\)，对有限片相加，得到 \(\rho\geq1\) 时所需结论。

最后，当 \(|\xi|\leq1\)，直接使用
\[
 |\widehat\mu(\xi)|
 \leq\int_\Sigma|a|\,\mathrm d\sigma<\infty.
\]
调整常数便得到整个频率空间上的估计。
\end{proof}

驻相系数还展示了曲面体积元为何必不可少。记 \(n=d-1\)，在图形片的驻点处，
\[
 |\det D_y^2\phi_\omega|
 =|\omega_d|^n\cdot|\det D^2h|
 =|\mathcal K|\cdot J^2.
\]
因此原图形振幅的几何部分满足
\[
 \frac{\eta(F(y))\cdot a(F(y))\cdot J(y)}
      {|\det D_y^2\phi_\omega(y)|^{1/2}}
 =\frac{\eta(F(y))\cdot a(F(y))}
      {|\mathcal K(y)|^{1/2}}.
\]
Jacobi 因子与 Hesse 行列式中的坐标因子恰好抵消。这与\secref{b2:sec:2102}中驻相首项的坐标不变性一致，也说明曲率下界趋于零时为何不能指望同一常数继续适用。

\ProseParagraphBreak

曲率退化必须按具体模型判断。若曲面是平面，沿法向的相位恒定，非零密度的 Fourier 变换可完全没有衰减。取整数 \(1\leq r<n\)，若局部相位恰为
\[
 \phi(u,v)=\frac{|u|^2}{2},
 \quad
 (u,v)\in\mathbb R^r\times\mathbb R^{n-r},
\]
且振幅是 \(b(u)c(v)\)，其中
\(b\in C_c^\infty(\mathbb R^r;\mathbb C)\)、
\(c\in C_c^\infty(\mathbb R^{n-r};\mathbb C)\)，则积分精确分解成
\[
 \left(\int_{\mathbb R^{n-r}}c(v)\,\mathrm dv\right)
 \left(\int_{\mathbb R^r}e^{\pi\mathrm i\lambda|u|^2}
                    \cdot b(u)\,\mathrm du\right).
\]
当 \(b(0)\int c\ne0\)，二次驻相公式给出非零的
\(\lambda^{-r/2}\) 主项。这个秩模型中，只有 \(r\) 个变量产生二次振荡；其余变量确实完全平直。圆柱曲面沿适当法向的局部模型具有这种性质。

\ProseParagraphBreak

不能把上述结论推广为“驻点 Hesse 矩阵的秩总能决定衰减”。例如对 \(m\geq2\)，相位
\[
 \phi(u,v)=\frac{|u|^2}{2}+v^{2m},
 \quad u\in\mathbb R^r,\quad v\in\mathbb R,
\]
在原点的 Hesse 矩阵仍只有秩 \(r\)，但最后一个变量并非平直。对乘积振幅，二次估计和 \(2m\) 阶 van der Corput 引理给出
\(O(\lambda^{-r/2-1/(2m)})\)。
同一个点上的二阶秩没有记录这个额外增益。一般退化曲率问题还要研究高阶项、驻点集合形状及频率方向；只有具备明确的法向非退化结构时，才可应用已证明的临界子流形公式。

\subsection{局部化}
\label{b2:subsec:210304}

局部化的目的不是简单地把积分切小，而是让所需下界和导数界在每一片中同时成立。若 \(\chi_j\) 是振幅支撑附近的有限光滑分割，则
\[
 I_\lambda(a,\phi)
 =\sum_j I_\lambda(a\chi_j,\phi).
\]
在非驻相片中记录梯度下界，在驻相片中记录 Hesse 矩阵的最小奇异值下界；分割函数导数进入新振幅。有限个固定片可以相加，但参数变化时，还必须检查能否保持同一覆盖和这些下界。曲面命题中的 \(K\times S^{d-1}\) 紧性，正是使方向分割和全部估计常数能够统一的原因。

\ProseParagraphBreak

一个方向逐渐变平的二次模型可以把这种依赖完全算清。取
\[
 J_{\lambda,\varepsilon}
 =\iint_{\mathbb R^2}
 e^{\pi\mathrm i\lambda(x^2+\varepsilon y^2)}
       \cdot a(x)b(y)\,\mathrm dx\,\mathrm dy,
 \quad \lambda\geq1,\quad\varepsilon\geq0,
\]
其中 \(a,b\in C_c^\infty(\mathbb R;\mathbb C)\) 固定。令
\[
 A(\lambda)=\int_{\mathbb R}e^{\pi\mathrm i\lambda x^2}
                           \cdot a(x)\,\mathrm dx,
 \quad
 B(\tau)=\int_{\mathbb R}e^{\pi\mathrm i\tau y^2}
                           \cdot b(y)\,\mathrm dy.
\]
Fubini 定理给出 \(J_{\lambda,\varepsilon}=A(\lambda)B(\lambda\varepsilon)\)。
对 \(A\) 用二阶 van der Corput，对 \(B\) 同时用普通积分界和二阶估计，可得统一上界
\[
 |J_{\lambda,\varepsilon}|
 \leq C_{a,b}\lambda^{-1/2}
 \begin{cases}
  1,&0\leq\lambda\varepsilon\leq1,\\
  (\lambda\varepsilon)^{-1/2},&\lambda\varepsilon\geq1.
 \end{cases}
\]
这条估计在 \(\varepsilon=0\) 也有意义。固定 \(\varepsilon>0\) 时，两个二次方向最终都产生衰减，得到
\(O(\lambda^{-1}\varepsilon^{-1/2})\)；
若 \(\lambda\varepsilon\) 很小，第二个方向尚未形成有效振荡，只能保留第一方向的 \(\lambda^{-1/2}\)。

\ProseParagraphBreak

若进一步取 \(a(0)\ne0\)、\(\int b\ne0\)，并令 \(\varepsilon=\lambda^{-2}\)，则
\[
 \lambda^{1/2}J_{\lambda,\lambda^{-2}}
 \longrightarrow
 e^{\pi\mathrm i/4}a(0)\int_{\mathbb R}b(y)\,\mathrm dy\ne0.
\]
这里第一因子用二次驻相首项，第二因子用 \(\tau=\lambda^{-1}\to0\) 时的控制收敛。因此即使每个正 \(\varepsilon\) 都给出可逆 Hesse 矩阵，也不存在对这些参数统一的 \(C\lambda^{-1}\) 界。这个例子直接检验了常数的退化，而没有改变驻点数目。后面的振荡积分算子中，混合 Hesse 矩阵的下界也必须按同样原则与有效参数一起追踪。

% !TeX root = ../../Book2.tex
\section{振荡积分算子}
\label{b2:sec:2104}

标量积分的衰减来自给定振幅下的相位抵消。算子范数却要同时控制整个 \(L^2\) 单位球；输入可以依赖于高频参数，甚至抵消相位的一部分。要找到在所有输入下仍然存在的增益，需要比较两个输出点的相位。这个比较将标量非驻相估计与 Hilbert 空间的 \(TT^*\) 方法联系起来。

\subsection{算子核}
\label{b2:subsec:210401}

设整数 \(d\ge1\)，\(U\subset\mathbb R^d_x\times\mathbb R^d_y\) 开，实值相位 \(\Phi\in C^\infty(U)\)，振幅 \(a\in C_c^\infty(U)\)。对 \(\lambda\ge1\) 定义
\[
 T_\lambda f(x)\coloneqq
 \int_{\mathbb R^d}e^{2\pi\mathrm i\lambda\Phi(x,y)}a(x,y)f(y)\,\mathrm dy.
\]
这里把振幅在 \(U\) 外延为零，只在其支撑附近使用相位。由于 \(y\) 支撑有界，Cauchy--Schwarz 不等式使积分对每个 \(x\) 都有意义；对 \(x\) 再积分，得到一个与 \(\lambda\) 无关的粗略 \(L^2\) 界。这个初步界保证伴随和算子复合已经定义，随后才研究高频衰减。

\ProseParagraphBreak
固定 \(x\) 后直接对核使用驻相法，会把 \(f\) 当作振幅的一部分，所得常数便涉及 \(f\) 的导数，而 \(L^2\) 单位球没有这种控制。更直接的障碍是分离相位
\(\Phi(x,y)=\psi(x)+\theta(y)\)。取
\(f_\lambda(y)=e^{-2\pi\mathrm i\lambda\theta(y)}g(y)\)，则
\[
 T_\lambda f_\lambda(x)
 =e^{2\pi\mathrm i\lambda\psi(x)}\int a(x,y)g(y)\,\mathrm dy.
\]
输入与输出上的指数因子都不改变 \(L^2\) 范数，因此一般不存在随 \(\lambda\) 衰减的算子界。

\ProseParagraphBreak
本节固定混合 Hesse 矩阵的行列顺序为
\[
 M(x,y)_{jk}\coloneqq
 \frac{\partial^2\Phi}{\partial y_j\partial x_k}(x,y),
 \quad 1\le j,k\le d.
\]
它把输出点的位移映成输入相位梯度的变化。分离相位的 \(M\) 恒为零，双线性相位 \(x\cdot y\) 的 \(M\) 则是单位矩阵。后一个模型在没有截断时满足
\(T_\lambda f(x)=\widehat f(-\lambda x)\)，Plancherel 给出准确的范数 \(\lambda^{-d/2}\)。混合矩阵的可逆性将使一般光滑相位在小片上具有同一衰减幂次。

\subsection{\texorpdfstring{\(TT^*\)}{TT*} 方法}
\label{b2:subsec:210402}

由伴随定义和 Fubini 定理，
\[
 T_\lambda^*g(y)=\int e^{-2\pi\mathrm i\lambda\Phi(x,y)}
                         \overline{a(x,y)}g(x)\,\mathrm dx.
\]
复合算子 \(T_\lambda T_\lambda^*\) 的核为
\[
 K_\lambda(x,z)=\int
 e^{2\pi\mathrm i\lambda(\Phi(x,y)-\Phi(z,y))}
 a(x,y)\overline{a(z,y)}\,\mathrm dy.
\]
这些换序可先对紧支撑有界输入进行。对于一般 \(L^2\) 输入，振幅各变量的紧支撑及 Cauchy--Schwarz 给出相同的绝对可积性，也可由稠密性延拓。相位差不再包含待估输入；因此非驻相法所要求的振幅导数，现在全部落在已知的 \(a\) 上。

\ProseParagraphBreak
Hilbert 空间中有
\[
 \|T_\lambda\|_{2\to2}^2
 =\|T_\lambda T_\lambda^*\|_{2\to2}.
\]
一边由算子范数的次乘性得到；另一边由
\(\|T_\lambda^*g\|_2^2
 =\langle T_\lambda T_\lambda^*g,g\rangle\)
及 \(\|T_\lambda^*\|=\|T_\lambda\|\) 得到。
在 \(x=z\) 上相位差为零，核不可能处处都小。有效的目标是在小片上证明它集中在宽度约 \(\lambda^{-1}\) 的对角邻域，再把这个邻域的体积转成算子范数。

\Needspace{6\baselineskip}
\begin{lemma}[Schur 检验]\label{b2:lem:21-schur}
设整数 \(m,n\ge1\)，\(K\) 是 \(\mathbb R^m\times\mathbb R^n\) 上的可测核，且存在有限常数 \(A,B\ge0\)，使
\[
 \mathop{\mathrm{ess\,sup}}_x\int_{\mathbb R^n}|K(x,z)|\,\mathrm dz\le A,
 \quad
 \mathop{\mathrm{ess\,sup}}_z\int_{\mathbb R^m}|K(x,z)|\,\mathrm dx\le B.
\]
则 \(Sg(x)=\int K(x,z)g(z)\,\mathrm dz\) 定义从 \(L^2(\mathbb R^n)\) 到 \(L^2(\mathbb R^m)\) 的有界算子，且
\(\|S\|_{2\to2}\le\sqrt{AB}\)。对每个 \(g\in L^2\)，该积分在几乎处处的 \(x\) 绝对收敛。
\end{lemma}
\begin{proof}
Tonelli 定理与第二个界给出
\[
 \int\!\int |K(x,z)|\cdot|g(z)|^2\,\mathrm dz\,\mathrm dx
 \le B\|g\|_2^2.
\]
因此内层积分几乎处处有限。在这些 \(x\) 上，带权 Cauchy--Schwarz 不等式给出绝对收敛和
\[
 |Sg(x)|^2
 \le\left(\int |K(x,z)|\,\mathrm dz\right)
       \left(\int |K(x,z)|\cdot|g(z)|^2\,\mathrm dz\right).
\]
对 \(x\) 积分并使用第一个界，得到
\(\|Sg\|_2^2\le AB\|g\|_2^2\)。
若最初只在有界紧支撑函数上定义 \(S\)，这个估计给出的连续延拓与上述几乎处处定义一致。
\end{proof}

\term[schur-jianyan]{Schur 检验}[Schur test]要求两个方向的积分界。对本节所需的核，若已经利用振荡证明
\[
 |K_\lambda(x,z)|\le C_N(1+\lambda|x-z|)^{-N},
\]
那么换元 \(w=\lambda(x-z)\) 给出
\[
 \int_{\mathbb R^d}(1+\lambda|x-z|)^{-N}\,\mathrm dz
 =\lambda^{-d}\int_{\mathbb R^d}(1+|w|)^{-N}\,\mathrm dw.
\]
取 \(N>d\)，右端有限，对另一个积分方向也一样。因此
\(\|T_\lambda T_\lambda^*\|_{2\to2}\le C\lambda^{-d}\)，开平方才得到 \(\lambda^{-d/2}\)。这里取核的绝对值是在非驻相估计之后；若一开始对原始振荡核取绝对值，只能回到不衰减的粗略界。

\subsection{\texorpdfstring{\(L^2\)}{L2} 估计}
\label{b2:subsec:210403}

\Needspace{6\baselineskip}
\begin{theorem}[非退化振荡算子的局部 \(L^2\) 估计]\label{b2:thm:21-oscillatory-l2}
设整数 \(d\ge1\)，\(U\subset\mathbb R^d_x\times\mathbb R^d_y\) 开，实值 \(\Phi\in C^\infty(U)\)，\((x_0,y_0)\in U\)，并设混合 Hesse 矩阵
\(M_{jk}=\partial_{y_j}\partial_{x_k}\Phi\) 在 \((x_0,y_0)\) 可逆。
则存在该点的乘积邻域 \(W\Subset U\)，使每个 \(a\in C_c^\infty(W)\) 所定义的算子
\[
 T_\lambda f(x)=\int e^{2\pi\mathrm i\lambda\Phi(x,y)}a(x,y)f(y)\,\mathrm dy
\]
满足
\[
 \|T_\lambda f\|_2\le C\lambda^{-d/2}\|f\|_2,
 \quad f\in L^2(\mathbb R^d),\quad\lambda\ge1.
\]
更一般地，若 \(a\in C_c^\infty(U)\)，且 \(M\) 在 \(\operatorname{supp}a\) 的邻域处处可逆，同一结论成立。常数可依赖于相位、振幅及其固定支撑，但与 \(f,\lambda\) 无关。
\end{theorem}
\begin{proof}
记 \(M_0=M(x_0,y_0)\)，其最小奇异值为 \(s_0>0\)。取有界凸邻域 \(X,Y\)，使
\(\overline X\times\overline Y\subset U\)，并且在其上
\(\|M(x,y)-M_0\|\le s_0/2\)。先设振幅支撑在 \(X\times Y\) 内。
当 \(x,z\in X\) 时，连接二者的线段仍在 \(X\) 中，故
\[
 \nabla_y\Phi(x,y)-\nabla_y\Phi(z,y)
 =\left[\int_0^1 M(z+t(x-z),y)\,\mathrm dt\right](x-z).
\]
括号内矩阵与 \(M_0\) 的距离不超过 \(s_0/2\)，所以相位梯度差的长度至少为 \((s_0/2)|x-z|\)。

对 \(x\ne z\) 定义归一化相位
\[
 \psi_{x,z}(y)\coloneqq
 \frac{\Phi(x,y)-\Phi(z,y)}{|x-z|}.
\]
它的梯度下界为 \(s_0/2\)。此外，对于每个多重指标 \(\beta\)，均值公式给出
\[
 \partial_y^\beta\psi_{x,z}(y)
 =\int_0^1\nabla_x\partial_y^\beta\Phi(z+t(x-z),y)
       \cdot\frac{x-z}{|x-z|}\,\mathrm dt.
\]
因此在固定紧集上，相位的所需导数一致有界，常数不随 \(x,z\) 接近而增大。振幅
\(a(x,y)\overline{a(z,y)}\) 的 \(y\) 导数及支撑体积也一致受控。
对参数 \(\mu=\lambda|x-z|\ge1\) 使用\cref{b2:prop:21-nonstationary-multi}，得到
\( |K_\lambda(x,z)|\le C_N\mu^{-N}\)。
当 \(\mu<1\) 或 \(x=z\) 时，直接用振幅的 \(L^1\) 上界。两种情形合并为
\[
 |K_\lambda(x,z)|\le C_N(1+\lambda|x-z|)^{-N}.
\]
在 \(X\) 外的输出点上核为零。取整数 \(N>d\)，由\cref{b2:lem:21-schur}及上一小节的换元公式得到 \(TT^*\) 的 \(C\lambda^{-d}\) 界，再开平方即得局部结论。

对于一般紧支撑，在其上选有限个上述乘积邻域，用从属光滑分割写成
\(a=\sum_{\nu=1}^L a_\nu\)。对应算子为
\(T_\lambda=\sum_{\nu=1}^L T_{\lambda,\nu}\)，因而
\[
 \|T_\lambda\|_{2\to2}
 \le\sum_{\nu=1}^L\|T_{\lambda,\nu}\|_{2\to2}
 \le\left(\sum_{\nu=1}^L C_\nu\right)\lambda^{-d/2}.
\]
有限求和完成整体结论，不需要估计不同小片之间的交叉核。
\end{proof}

小片条件在证明中有明确用途：平均矩阵仍接近同一个可逆矩阵。仅知每一点的混合矩阵可逆，不能推出它沿长线段的平均可逆，也不能推出 \(x\mapsto\nabla_y\Phi(x,y)\) 在大区域上单射。因此先分别估计局部算子，再求和。每个小片内的核集中在对应的对角线附近，整体核则可能具有更复杂的集中结构。

\subsection{参数依赖性}
\label{b2:subsec:210404}

在参数族中，要统一上述证明，需要同时控制支撑、小片半径、逆矩阵和分部积分系数。一个充分条件是：所有振幅支撑在同一紧集 \(K\Subset U\)，相位在 \(K\) 的同一邻域上定义，混合矩阵最小奇异值一致不小于 \(c>0\)，并且所需的相位、振幅导数一致有界。相位三阶导数界使混合矩阵有统一的连续模，因而可选择统一的小片半径；\(K\) 的紧性再给出片数及分割函数导数的统一控制。

\ProseParagraphBreak
对于一个给定的 \(L^2\) 估计，只需有限阶光滑性。固定整数 \(N>d\) 后，在上述固定邻域上控制相位到 \(N+2\) 阶、振幅到 \(N\) 阶，是保证这套证明中常数统一的一个充分条件。相位的这一阶数还包含均值公式所多用的一次 \(x\) 微分。若要求任意阶渐近展开，则要对每一阶分别提出相应导数界，不能从一次有限阶估计推出全部展开。

\ProseParagraphBreak
仅记录行列式不为零，尚不足以核对这套证明的统一常数。例如
\(M_\epsilon=\operatorname{diag}(\epsilon,\epsilon^{-1})\)
的行列式恒为一，但最小奇异值趋零，矩阵范数趋于无穷。这个例子表明上述充分条件失效，却不单独证明对应算子范数变坏：无截断的线性 Fourier 模型，其准确范数由行列式决定。若要证明某族的统一衰减确实失败，应直接检验算子。

\ProseParagraphBreak
一个直接检验是 \(\Phi_\epsilon(x,y)=\epsilon x\cdot y\)，其中 \(0<\epsilon\le1\)。取固定分离振幅 \(\chi(x)\eta(y)\)，则
\[
 T_{\lambda,\epsilon}f(x)
 =\chi(x)\widehat{\eta f}(-\lambda\epsilon x).
\]
当 \(\lambda\epsilon\ge1\) 时，Plancherel 给出的上界为
\(C(\lambda\epsilon)^{-d/2}\|f\|_2\)；当 \(\lambda\epsilon<1\) 时，紧支撑给出 \(C\|f\|_2\)。这两个界说明真正控制振荡的是 \(\lambda\epsilon\)。若选择固定 \(g\) 使 \(\int\eta g\ne0\)，并令 \(\epsilon=\lambda^{-2}\)，则
\[
 T_{\lambda,\epsilon}g\longrightarrow
 \chi\int\eta(y)g(y)\,\mathrm dy
 \quad\text{于 }L^2.
\]
只要 \(\chi\not\equiv0\)，极限非零，故不可能有独立于 \(\epsilon\) 的 \(C\lambda^{-d/2}\) 界。

\ProseParagraphBreak
在固定非退化情形，\(\lambda^{-d/2}\) 本身也不能改善。仍取双线性相位和分离振幅，设 \(\chi(0)\ne0\)，固定 \(g\in C_c^\infty\) 满足 \(\eta g\not\equiv0\)。则
\[
 \|T_\lambda g\|_2^2
 =\lambda^{-d}\int |\chi(v/\lambda)|^2\cdot
                     |\widehat{\eta g}(-v)|^2\,\mathrm dv.
\]
支配收敛及 Plancherel 给出
\(\lambda^d\|T_\lambda g\|_2^2
 \to |\chi(0)|^2\cdot\|\eta g\|_2^2>0\)，所以算子范数至少是 \(c\lambda^{-d/2}\)。这个测试无需让输入依赖于 \(\lambda\)。

\ProseParagraphBreak
若输入与输出变量的维数不同，混合矩阵不再是方阵，应改为考虑其秩及无振荡变量的作用。多加一个仅作普通积分的变量，并不会多产生半个幂次的衰减；衰减来自真正配对的振荡方向。后续限制问题中出现的秩与曲率条件，需要在相应几何下重新分析。

% !TeX root = ../../Book2.tex
\OptionalSection{Fourier 积分算子初步}{b2:sec:2105}

前几节从相位的导数得到积分或算子范数估计。本节进一步问：输入在某个位置、某个频率方向的振荡，会传到输出的哪里？Fourier 积分算子的相位可以同时记录这两类信息。从驻点关系中可以读出输入和输出的位置、频率，再考察改变相位参数时这些信息怎样保持。相关背景可参阅\appref{b2:app:G}及专著 \cite{Hrmander2009AnalysisIV}。

\subsection{非退化相位}
\label{b2:subsec:210501}

先考虑不会产生收敛困难的表示。设 \(U\subset\mathbb R^d\) 开，实值 \(\phi\) 光滑，\(a\in C_c^\infty(U\times\mathbb R^d)\)，\(f\in\mathcal S(\mathbb R^d)\)。由 Fourier 变换的定义及 Fubini 定理，
\[
 \begin{aligned}
 Af(x)
 &=\iint e^{2\pi\mathrm i(\phi(x,\theta)-y\cdot\theta)}
          a(x,\theta)f(y)\,\mathrm dy\,\mathrm d\theta\\
 &=\int e^{2\pi\mathrm i\phi(x,\theta)}
          a(x,\theta)\widehat f(\theta)\,\mathrm d\theta.
 \end{aligned}
\]
频率支撑紧时，对 \(x,y\) 作任意次微分都可以移入核积分，所以核光滑。这个情形适合先检查相位和换元，但一般 Fourier 积分算子的奇性来自无界频率，不能由紧频率核本身看出。

\ProseParagraphBreak
例如 \(\phi(x,\theta)=x\cdot\theta\) 给出左量子化的伪微分表示。取 \(a\equiv1\) 后，积分对 Schwartz 输入仍绝对收敛，由 Fourier 反演得到恒等算子。若用 \(\chi(\theta/R)\) 截断，其中 \(\chi\in C_c^\infty\) 在零点附近等于一，则各个截断核都光滑，而它们只在分布意义下趋于 \(\delta(x-y)\)。这是光滑核序列产生奇异极限的一个基本例子。

\ProseParagraphBreak
另一个模型是
\[
 \phi_t(x,\theta)=x\cdot\theta+t|\theta|,
 \quad
 A_tf=\mathcal F^{-1}\bigl(e^{2\pi\mathrm i t|\theta|}\widehat f(\theta)\bigr).
\]
这个 Fourier 乘子定义了半波算子 \(e^{2\pi\mathrm i t|D|}\)。它的乘子模为一，故由 Plancherel 延拓为 \(L^2\) 上的酉算子。其相位只在 \(\theta\ne0\) 光滑；低频部分将另行处理。以下的齐次相位描述均针对非零频率。

\Needspace{6\baselineskip}
\begin{definition}[齐次非退化相位]\label{b2:def:21-nondegenerate-phase}
设整数 \(d\ge1\)，\(U,V\subset\mathbb R^d\) 开，\(\Gamma\subset\mathbb R^k\setminus\{0\}\) 为开锥，\(k\ge1\)。实值光滑函数 \(\Psi(x,y,\theta)\) 若满足
\(\Psi(x,y,r\theta)=r\Psi(x,y,\theta)\)（\(r>0\)），且在临界集合
\[
 C_\Psi\coloneqq\{(x,y,\theta):\partial_\theta\Psi(x,y,\theta)=0\}
\]
上，\(k\) 个微分
\(\mathrm d_{x,y,\theta}(\partial_{\theta_j}\Psi)\)
线性无关，则称为\term[qi-ci-fei-tui-hua-xiang-wei]{齐次非退化相位}[homogeneous nondegenerate phase]。本节只在
\(\partial_x\Psi\ne0\)、\(\partial_y\Psi\ne0\)
的部分讨论输入、输出均为非零的余切方向。
\end{definition}

这里的非退化是驻点方程作为约束的正则性。隐函数定理说明 \(C_\Psi\) 的维数为 \(2d\)，因为总变量数为 \(2d+k\)，有 \(k\) 个独立约束。它并不要求频率 Hesse 矩阵\(\partial_\theta^2\Psi\)可逆。事实上，Euler 齐次恒等式及其微分给出
\[
 \Psi=\theta\cdot\partial_\theta\Psi,
 \quad (\partial_\theta^2\Psi)\theta=0;
\]
因而频率 Hesse 矩阵总有径向零方向。

\ProseParagraphBreak
对于混合表示 \(\Psi=\phi(x,\theta)-y\cdot\theta\)，约束微分中的 \(y\) 分块是 \(-I\)，所以这组约束自动独立。若再要求 \(\partial_x\partial_\theta\phi\) 可逆，才可从
\(y=\partial_\theta\phi(x,\theta)\) 局部解出 \(x\)，将所得关系写成输入到输出的局部图。前几节还有另一种条件：标量驻相要求积分变量上的二阶 Hesse 矩阵可逆。三种条件分别服务于驻相渐近、\(L^2\) 算子估计和驻点集合的几何，使用时必须说明求导变量。

\subsection{典范关系}
\label{b2:subsec:210502}

\Needspace{6\baselineskip}
\begin{definition}[相位所生成的典范关系]\label{b2:def:21-canonical-relation}
在\cref{b2:def:21-nondegenerate-phase}的条件下，相位 \(\Psi\) 所生成的\term[dianfan-guanxi]{典范关系}[canonical relation]是集合
\[
 \mathcal C_\Psi\coloneqq
 \bigl\{(x,\partial_x\Psi;\ y,-\partial_y\Psi):
       (x,y,\theta)\in C_\Psi\bigr\},
\]
并只保留两端余切向量均非零的部分。这里 \((x,\xi)\) 表示输出位置及频率，\((y,\eta)\) 表示输入位置及频率。
\end{definition}

输入端的负号来自核的约定。若输入局部含有正频率振荡
\(e^{2\pi\mathrm i y\cdot\eta}\)，则总相位对 \(y\) 的驻点条件为
\(\partial_y\Psi+\eta=0\)，而输出相位梯度是 \(\partial_x\Psi\)。因此上式记录的正是可能相互耦合的输入和输出频率。

\ProseParagraphBreak
非退化条件也保证上述参数化在局部没有额外的切向冗余。若 \(C_\Psi\) 上一个切向量在 \((x,\partial_x\Psi;y,-\partial_y\Psi)\) 下的像为零，则其 \(\delta x,\delta y\) 均为零，并且
\[
 \Psi_{x\theta}\delta\theta=0,
 \quad\Psi_{y\theta}\delta\theta=0,
 \quad\Psi_{\theta\theta}\delta\theta=0.
\]
最后一个等式来自临界约束的微分。把三个式子合并，就是约束微分矩阵的转置作用于 \(\delta\theta\) 为零；其行满秩迫使 \(\delta\theta=0\)。这说明关系具有正确的局部维数，但不排除不同参数点给出同一像。Euler 恒等式又使 \(\Psi\) 在 \(C_\Psi\) 上等于零，从而沿该集合有
\(\partial_x\Psi\cdot\mathrm dx+\partial_y\Psi\cdot\mathrm dy=0\)。这些恒等式将相位的微分与输入、输出频率联系起来，是进一步研究其几何结构的出发点。

\ProseParagraphBreak
恒等相位 \(\Psi=(x-y)\cdot\theta\) 的约束为 \(x=y\)，所以
\[
 \mathcal C_\Psi=\{(x,\theta;x,\theta):\theta\ne0\}.
\]
半波相位则给出
\[
 \Psi_t=(x-y)\cdot\theta+t|\theta|,
 \quad x-y+t\frac{\theta}{|\theta|}=0,
\]
于是
\[
 \mathcal C_{\Psi_t}
 =\left\{\left(y-t\frac{\theta}{|\theta|},\theta;\ y,\theta\right):
          \theta\ne0\right\}.
\]
对本节采用的正指数 \(e^{2\pi\mathrm i t|D|}\)，输出位置为
\(x=y-t\theta/|\theta|\)；若改成负指数，方向反转。位置与频率必须同时记录，单凭空间中的球面无法恢复这个符号信息。

\subsection{相位等价}
\label{b2:subsec:210503}

首先看真正的频率坐标变换。设
\(\theta=\Theta(x,y,\eta)\) 对每个 \((x,y)\) 都是两个开锥之间的光滑微分同胚，并且对 \(\eta\) 正齐次一次。换元后的相位及振幅为
\[
 \widetilde\Psi(x,y,\eta)=\Psi(x,y,\Theta(x,y,\eta)),
 \quad
 \widetilde a=a(x,y,\Theta(x,y,\eta))
              |\det D_\eta\Theta|.
\]
对于紧频率支撑，这是普通换元公式。其驻点方程为
\[
 \partial_\eta\widetilde\Psi
 =(D_\eta\Theta)^{\mathsf T}\partial_\theta\Psi=0,
\]
与原约束等价。在临界集合上，对 \(x,y\) 微分时由 \(\Theta\) 产生的额外项均乘着 \(\partial_\theta\Psi=0\)，故两端余切向量相同。约束的微分也只多乘可逆矩阵及作坐标变换，因此非退化性与典范关系都保持不变。Jacobi 因子则必须保留，否则积分算子会改变。

\ProseParagraphBreak
增加辅助二次变量也能产生等价表示，但需要顾及齐次性。直接加 \(|\eta|^2/2\) 会破坏一次齐次；在 \(\theta\ne0\) 上应改用
\[
 \widetilde\Psi(x,y,\theta,\eta)
 =\Psi(x,y,\theta)+\frac{|\eta|^2}{2|\theta|},
 \quad\eta\in\mathbb R^m.
\]
它对 \((\theta,\eta)\) 同时伸缩为一次齐次。新的约束先给出 \(\eta=0\)，再给出旧约束 \(\partial_\theta\Psi=0\)，所以典范关系相同；新约束的 \(\eta\) 分块为 \(|\theta|^{-1}I\)，保证其独立性。

\ProseParagraphBreak
为核对算子而不只核对驻点集合，可先把 \(\theta\) 限制在远离零的紧集。\cref{b2:prop:21-complex-gaussian}逐坐标给出 Abel 正则化积分
\[
 \lim_{\epsilon\downarrow0}\int_{\mathbb R^m}
 e^{\pi\mathrm i|\eta|^2/|\theta|}
 e^{-\pi\epsilon|\eta|^2}\,\mathrm d\eta
 =e^{\mathrm i\pi m/4}|\theta|^{m/2}.
\]
因此把振幅换成
\(e^{-\mathrm i\pi m/4}|\theta|^{-m/2}a(x,y,\theta)\)，再加入这组辅助变量，正则化后恰恢复原核。该等式在固定的紧频率集上一致，允许把极限移入 \(\theta\) 积分。对于无界频率，因子 \(|\theta|^{-m/2}\) 还会改变振幅的阶；必须连同符号条件及分布极限一起处理。

\subsection{微局部观点}
\label{b2:subsec:210504}

微局部观察先用空间截断选定位置，再用锥形频率截断选定方向。相位在该方向没有驻点时，可以通过反复分部积分排除相应核奇性。为了说明这句话在无界频率中为什么成立，需要把每个有限频率壳层的估计求和，而不是直接把紧支撑非驻相公式套在整个频率空间上。

\ProseParagraphBreak
具体地，在一个固定的 \((x,y)\) 紧邻域及频率开锥内，设 \(\Psi\) 是一次齐次光滑相位，并且在紧包含于该开锥的单位球面截面内的角向支撑上
\( |\partial_\theta\Psi|\ge c>0\)。设高频振幅 \(a\) 满足：对所有多重指标 \(\alpha,\beta\)，存在常数使
\[
 |\partial_{x,y}^\alpha\partial_\theta^\beta a(x,y,\theta)|
 \le C_{\alpha\beta}(1+|\theta|)^{m-|\beta|},
 \quad |\theta|\ge1.
\]
把频率分为 \(|\theta|\asymp2^j\) 的光滑壳层，令 \(\theta=2^j\omega\)。壳层核变成固定环域上的振荡积分，其参数为 \(2^j\)，Jacobi 因子为 \(2^{jk}\)，振幅规模至多为 \(2^{jm}\)。对 \((x,y)\) 作 \(|\alpha|\) 次微分，最坏再增加 \(2^{j|\alpha|}\)。利用非驻相估计任意 \(N\) 次，便得到
\[
 |\partial_{x,y}^\alpha K_j(x,y)|
 \le C_{\alpha,N}\,2^{j(k+m+|\alpha|-N)}.
\]
取 \(N>k+m+|\alpha|\)，右端可求和。对每个微分阶数分别如此选择 \(N\)，壳层级数及其全部导数一致收敛，所以这个位置、方向片上的核光滑。

\ProseParagraphBreak
在恒等算子模型，\(\partial_\theta\Psi=x-y\)，上述论证解释了核在对角线外光滑。在半波模型，条件变为
\(x-y+t\theta/|\theta|\ne0\)，因此高频贡献只能在满足传播关系的方向上产生核奇性。存在驻点只是可能产生奇性的条件；振幅若在相关位置消失，相应贡献仍可能消失。一般的波前集映射需要进一步的局部化论证，不能只由这一单向判别替代。

\ProseParagraphBreak
半波的低频核可直接处理。若 \(\chi\) 紧支撑于零点附近，则
\[
 K_{t,\mathrm{low}}(x,y)
 =\int e^{2\pi\mathrm i((x-y)\cdot\theta+t|\theta|)}
          \chi(\theta)\,\mathrm d\theta
\]
对 \(x,y\) 的任何微分只引入频率多项式，仍是绝对可积函数，故核光滑。因频率支撑有界，所有频率矩都可积，空间变量的各阶导数因而都有绝对收敛的积分表示。分离低频之后，零点的不光滑与高频传播的几何便不会混淆。完整波前集定义、伪微分椭圆正则性与传播定理见\appref{b2:app:G}，其中分别说明各结论的证明范围。


\begin{chaptersummary}
非驻相法把相位导数下界转为反复分部积分；端点不消失时，留下的边界项本身就可能是主项。van der Corput 引理通过小导数区域的测度控制处理有限阶退化。非退化驻相则给出更细的展开：每个有效二次方向贡献半个衰减幂次，行列式给出大小，惯性指数给出复相位，坐标换元的体积因子进入振幅。

\ProseParagraphBreak

估计应连同所控制的对象一起使用。固定振幅的导数界、随参数缩小的支撑，以及允许输入自由变化的算子单位球，不能彼此替代。振荡算子的混合 Hesse 矩阵控制相位梯度在两个输出点之间的变化；\(TT^*\) 在各个小片上先得到一个集聚在对角线附近的核，再由 Schur 估计和开平方给出 \(L^2\) 增益。有限个小片的估计相加，便得到紧支集振幅上的整体界。

\ProseParagraphBreak

遇到参数退化，先把参数与空间尺度共同规范化，再判断哪一种估计仍然一致。振幅的消失、部分方向的非退化和不同驻点之间的干涉，都可能改变实际衰减。选读部分进一步用位置与频率方向描述传播；它提供几何信息，而具体算子的大小与误差仍要由相应估计验证。
\end{chaptersummary}

\BookExercises

% 自拟；基础到中等：由端点消失阶求非线性非驻相边界首项，并验证单端贡献的尖锐阶。
\Needspace{6\baselineskip}
\begin{exercise}\label{b2:ex:21-new-1}
设整数 \(m\geq0\)，\(a\in C^\infty([0,1];\mathbb C)\)、\(\phi\in C^\infty([0,1];\mathbb R)\)，且 \(|\phi'|\geq c>0\)。假定
\(a^{(j)}(0)=a^{(j)}(1)=0\) 对 \(0\leq j<m\) 成立；\(m=0\) 时这一条件为空。证明当 \(\lambda\to+\infty\) 时，
\[
 \begin{aligned}
 I_\lambda
 &=\int_0^1 e^{2\pi\mathrm i\lambda\phi(t)}\cdot a(t)\,\mathrm dt\\
 &=\frac{(-1)^m}{(2\pi\mathrm i\lambda)^{m+1}}
 \left[\frac{e^{2\pi\mathrm i\lambda\phi(t)}\cdot a^{(m)}(t)}
                  {\phi'(t)^{m+1}}\right]_{t=0}^{t=1}
 +O(\lambda^{-m-2}).
 \end{aligned}
\]
取 \(\phi(t)=t+t^2\)、\(a(t)=t^m(1-t)^{m+1}\)，求出非零首项，并说明由此为何不能把一般结论改为 \(o(\lambda^{-m-1})\)。
\end{exercise}
\begin{solution}
记 \(\kappa=2\pi\mathrm i\lambda\)，并递归令
\(b_0=a\)、\(b_{j+1}=(b_j/\phi')'\)。每次分部积分给出
\[
 I_\lambda
 =\sum_{j=0}^{M-1}(-1)^j\kappa^{-j-1}
       \left[\frac{e^{\kappa\phi}\cdot b_j}{\phi'}\right]_0^1
   +(-1)^M\kappa^{-M}
       \int_0^1e^{\kappa\phi}\cdot b_M\,\mathrm dt.
\]
归纳展开递推式，\(b_j\) 的最高阶振幅导数项为
\((\phi')^{-j}\cdot a^{(j)}\)，其余项只含低于 \(j\) 阶的振幅导数，系数由 \(\phi\) 的导数及 \((\phi')^{-1}\) 构成。因此在两个端点，\(j<m\) 时有 \(b_j=0\)，而
\(b_m=(\phi')^{-m}\cdot a^{(m)}\)。取 \(M=m+2\)，第 \(m+1\) 个边界项和最后的积分都是 \(O(\lambda^{-m-2})\)；它们的系数有限，因为区间紧且 \(|\phi'|\geq c\)。剩余第 \(m\) 个边界项就是所求公式。

具体例中，\(a^{(m)}(0)=m!\)、\(a^{(m)}(1)=0\)、\(\phi(0)=0\)、\(\phi'(0)=1\)，故
\[
 I_\lambda
 =\frac{(-1)^{m+1}m!}{(2\pi\mathrm i\lambda)^{m+1}}
       +O(\lambda^{-m-2}),
 \quad
 \lambda^{m+1}|I_\lambda|
       \longrightarrow\frac{m!}{(2\pi)^{m+1}}>0.
\]
这个例子只有左端贡献非零，所以没有两端首项相互抵消的问题。一般情况下，两个端点即使都贡献非零系数，它们的和也可能沿某些参数抵消；因此总和在某些大参数处可能小于各首项的量级。
\end{solution}

% 自拟；中等：退化单侧相位配消失振幅，证明振荡常数非零并控制无穷远尾项，不仅作形式缩放。
\Needspace{6\baselineskip}
\begin{exercise}\label{b2:ex:21-new-2}
设整数 \(k\geq2\)、\(0\leq m\leq k-2\)，\(\chi\in C_c^\infty([0,\infty);\mathbb R)\) 在零点附近恒为一；这里光滑指能够延拓到零点的一个邻域。令
\[
 I_{k,m}(\lambda)=\int_0^\infty
       e^{2\pi\mathrm i\lambda t^k}\cdot t^m\cdot\chi(t)\,\mathrm dt,
 \quad \alpha=\frac{m+1}{k}.
\]
证明反常积分
\[
 C_{k,m}=\int_0^\infty e^{2\pi\mathrm i s^k}\cdot s^m\,\mathrm ds
\]
收敛且虚部严格为正，并且对每个整数 \(N\geq1\)，
\[
 I_{k,m}(\lambda)
 =C_{k,m}\lambda^{-(m+1)/k}+O_N(\lambda^{-N}).
\]
可在替换 \(v=2\pi s^k\) 后，将正弦积分的相邻正、负半周期配对。另求 \(m=k-1\) 时的首项，并解释这一端点为什么不再由上面普通收敛的反常积分给出。
\end{exercise}
\begin{solution}
换元后，
\[
 C_{k,m}=\frac1{k(2\pi)^\alpha}
          \int_0^\infty e^{\mathrm iv}\cdot v^{\alpha-1}\,\mathrm dv,
 \quad 0<\alpha<1.
\]
零点附近绝对可积。在无穷远对 \(e^{\mathrm iv}\) 积分分部，边界因子 \(v^{\alpha-1}\) 趋零，导数 \((\alpha-1)v^{\alpha-2}\) 绝对可积，故反常积分收敛。其正弦部分在第 \(j\) 对半周期上的积分为
\[
 \int_0^\pi\sin v\cdot
 \left[(2j\pi+v)^{\alpha-1}
       -\bigl((2j+1)\pi+v\bigr)^{\alpha-1}\right]\,\mathrm dv>0
 \quad(j\geq0).
\]
方括号为正，因为幂函数严格递减；第一项在 \(j=0\) 时仍可积。按整周期取部分积分再取极限，每一项为正且第一项非零，所以总虚部严格为正。

在完整的反常积分中作 \(s=\lambda^{1/k}t\)，得到
\[
 \int_0^\infty e^{2\pi\mathrm i\lambda t^k}\cdot t^m\,\mathrm dt
       =\lambda^{-\alpha}C_{k,m}.
\]
因此两者之差是振幅 \(q_0(t)=t^m\cdot\bigl(\chi(t)-1\bigr)\) 的振荡积分。这个振幅在零点附近消失，在足够大的 \(t\) 上等于 \(-t^m\)。令
\[
 q_{j+1}(t)=\left(\frac{q_j(t)}{kt^{k-1}}\right)'.
\]
所有 \(q_j\) 在零点附近都为零；在无穷远有
\(q_j(t)=O(t^{m-jk})\)。分部积分的边界因子
\(q_j(t)/(kt^{k-1})=O(t^{m-(j+1)k+1})\)
趋零，而 \(q_N\in L^1\) 对 \(N\geq1\) 成立。先在有限区间积分分部，再令上端趋于无穷，因而得到
\[
 I_{k,m}(\lambda)-\lambda^{-\alpha}C_{k,m}
 =(-2\pi\mathrm i\lambda)^{-N}
       \int_0^\infty e^{2\pi\mathrm i\lambda t^k}\cdot q_N(t)\,\mathrm dt.
\]
取 \(L^1\) 上界即得任意阶余项。由于 \(C_{k,m}\ne0\)，还得到
\(\lambda^{(m+1)/k}|I_{k,m}(\lambda)|\to|C_{k,m}|\)，所以这里的幂次是实际大小。

当 \(m=k-1\) 时，令 \(u=t^k\)，有
\[
 I_{k,k-1}(\lambda)
 =\frac1k\int_0^\infty e^{2\pi\mathrm i\lambda u}
                         \cdot\chi(u^{1/k})\,\mathrm du.
\]
新振幅在零点附近为一，其余部分光滑且紧支撑；它的导数远离零点。一次分部积分留下左端项，余项可任意次积分分部，故
\[
 I_{k,k-1}(\lambda)
 =-\frac1{2\pi\mathrm i k\lambda}+O_N(\lambda^{-N}).
\]
此时未经截断的积分变成常数倍 \(\int_0^\infty e^{\mathrm iv}\,\mathrm dv\)，普通反常极限不存在。非零首项来自保留下来的端点，而非上一个收敛论证。
\end{solution}

% 自拟；中等：由凸相位的参数临界值导出逆Hesse矩阵，再将驻相首项与梯度换元、Plancherel的能量配平。
\Needspace{6\baselineskip}
\begin{exercise}\label{b2:ex:21-new-3}
设整数 \(d\geq1\)，\(\psi\in C^\infty(\mathbb R^d;\mathbb R)\)，且
\(D^2\psi(x)\geq cI\) 对所有 \(x\) 成立，其中 \(c>0\)。设 \(a\in C_c^\infty(\mathbb R^d;\mathbb C)\)，对 \(\lambda\geq1\) 令
\[
 I_\lambda(p)=\int_{\mathbb R^d}
   e^{2\pi\mathrm i\lambda(\psi(x)-p\cdot x)}\cdot a(x)\,\mathrm dx.
\]
证明每个 \(p\in\mathbb R^d\) 都对应唯一的 \(x(p)\)，使 \(\nabla\psi\bigl(x(p)\bigr)=p\)。令
\(S(p)=\psi\bigl(x(p)\bigr)-p\cdot x(p)\)，求 \(\nabla S\) 与 \(D^2S\)，并求 \(I_\lambda(p)\) 的驻相首项；要求余项在每个紧参数集上一致。

\ProseParagraphBreak

若将首项写为
\(e^{2\pi\mathrm i\lambda S(p)}\lambda^{-d/2}e^{\mathrm i\pi d/4}\cdot b(p)\)，证明
\[
 \int_{\mathbb R^d}|b(p)|^2\,\mathrm dp
       =\int_{\mathbb R^d}|a(x)|^2\,\mathrm dx.
\]
再直接用 Plancherel 求 \(\int|I_\lambda(p)|^2\,\mathrm dp\)，解释它与上述恒等式为何相符。
\end{exercise}
\begin{solution}
二阶 Taylor 积分公式给出
\[
 \psi(x)-p\cdot x
 \geq\psi(0)+(\nabla\psi(0)-p)\cdot x+\frac c2|x|^2.
\]
因此它在无穷远趋于正无穷，连续性使最小值在某个紧球内达到。最小点满足 \(\nabla\psi(x)=p\)。对任意 \(x,z\)，沿线段积分又有
\[
 (\nabla\psi(x)-\nabla\psi(z))\cdot(x-z)
       \geq c|x-z|^2,
\]
故梯度映射单射。它处处导数可逆，逆函数定理说明全局逆映射 \(x(p)\) 光滑，并满足
\[
 D_px(p)=\bigl[D^2\psi\bigl(x(p)\bigr)\bigr]^{-1}.
\]
对 \(S\) 求导时，来自 \(\psi\bigl(x(p)\bigr)\) 的链式项与 \(p\cdot x(p)\) 中的对应项相消，所以
\[
 \nabla S(p)=-x(p),
 \quad D^2S(p)=-\bigl[D^2\psi\bigl(x(p)\bigr)\bigr]^{-1}.
\]
这个符号取决于题中选取 \(\psi(x)-p\cdot x\) 的次序。

相位有唯一驻点 \(x(p)\)，其 Hesse 矩阵正定，故\cref{b2:thm:21-stationary-phase}给出
\[
 I_\lambda(p)=e^{2\pi\mathrm i\lambda S(p)}\lambda^{-d/2}
 e^{\mathrm i\pi d/4}\cdot
 \frac{a\bigl(x(p)\bigr)}{\det\Bigl(D^2\psi\bigl(x(p)\bigr)\Bigr)^{1/2}}
 +O(\lambda^{-d/2-1}).
\]
给定紧参数集 \(P\)，先取包含它且闭包紧的参数开邻域 \(V\)，再取一个紧集 \(K\)，包含 \(\operatorname{supp}a\) 和所有 \(p\in\overline V\) 对应的驻点。这样的 \(K\) 存在，因为 \(x(p)\) 连续。每个 \(p\in V\) 的驻点都在同一个 \(K\) 内，Hesse 矩阵有同一个正下界，所需导数在相应紧邻域内也一致有界。因此\cref{b2:prop:21-parametric-stationary}可一致应用。若驻点不在振幅支撑内，首项为零；围绕该驻点局部化后，余下积分仍由统一非驻相估计处理。

梯度映射是全局微分同胚，Jacobi 行列式为正数 \(\det D^2\psi(x)\)。上述首项中的 \(b\) 因而满足
\[
 \int|b(p)|^2\,\mathrm dp
 =\int\frac{|a(x)|^2}{\det D^2\psi(x)}
                    \cdot\det D^2\psi(x)\,\mathrm dx
 =\|a\|_2^2.
\]
另一方面，令 \(g_\lambda=e^{2\pi\mathrm i\lambda\psi}\cdot a\)，则
\(I_\lambda(p)=\widehat g_\lambda(\lambda p)\)。Plancherel 与换元直接给出准确等式
\[
 \int|I_\lambda(p)|^2\,\mathrm dp
       =\lambda^{-d}\cdot\|g_\lambda\|_2^2
       =\lambda^{-d}\cdot\|a\|_2^2.
\]
驻相首项的平方积分恰有同一个值。梯度换元给出首项的平方积分，Plancherel 则给出原振荡积分的准确范数，两者由此相互对应。
\end{solution}

% 自拟；中等综合：固定有限维的分离双线性模型，逐坐标证明含退化参数的尖锐两侧范数。
\Needspace{6\baselineskip}
\begin{exercise}\label{b2:ex:21-new-4}
固定整数 \(d\geq1\) 及实函数 \(\chi\in C_c^\infty(\mathbb R)\)，满足 \(\chi\geq0\)，且在 \([-1,1]\) 上恒为一。设 \(\lambda\geq1\)、\(\varepsilon_1,\ldots,\varepsilon_d\geq0\)，对 \(f\in L^2(\mathbb R^d)\) 定义
\[
 T_{\lambda,\varepsilon}f(x)
 =\left(\prod_{j=1}^d\chi(x_j)\right)
   \int_{\mathbb R^d}
     e^{2\pi\mathrm i\lambda\sum_{j=1}^d\varepsilon_jx_jy_j}
     \cdot\left(\prod_{j=1}^d\chi(y_j)\right)\cdot f(y)\,\mathrm dy.
\]
令 \(w(\mu)=1\) 当 \(0\leq\mu\leq1\)，\(w(\mu)=\mu^{-1/2}\) 当 \(\mu>1\)。证明存在只依赖 \(d,\chi\) 的正数 \(c_d,C_d\)，使
\[
 c_d\prod_{j=1}^dw(\lambda\varepsilon_j)
 \leq\|T_{\lambda,\varepsilon}\|_{2\to2}
 \leq C_d\prod_{j=1}^dw(\lambda\varepsilon_j).
\]
先证明一维结论：上界分别用核的平方积分和缩放 Plancherel；下界可固定一个非负光滑输入，再利用其 Fourier 变换在原点附近不为零。最后用 Fubini 与乘积输入处理 \(d\) 维。
\end{exercise}
\begin{solution}
在一维记
\[
 S_\mu f(x)=\chi(x)\cdot
       \int e^{2\pi\mathrm i\mu xy}\cdot\chi(y)\cdot f(y)\,\mathrm dy.
\]
核的平方积分范数为 \(\|\chi\|_2^2\)，所以对所有 \(\mu\geq0\)，
\(\|S_\mu\|\leq\|\chi\|_2^2\)。当 \(\mu>0\) 时，
\(S_\mu f(x)=\chi(x)\cdot\widehat{\chi\cdot f}(-\mu x)\)，故
\[
 \|S_\mu f\|_2
 \leq\|\chi\|_\infty\cdot\mu^{-1/2}\cdot\|\chi\cdot f\|_2
 \leq\|\chi\|_\infty^2\cdot\mu^{-1/2}\cdot\|f\|_2.
\]
合并两种界得到 \(\|S_\mu\|\leq C\cdot w(\mu)\)。零参数时使用第一个界，不对零作负幂运算。

取固定的非零非负 \(g\in C_c^\infty\bigl((-1,1)\bigr)\)，归一化使 \(\|g\|_2=1\)。因为 \(\widehat g(0)=\int g>0\)，连续性给出 \(\delta\in(0,1]\)、\(c_0>0\)，使 \(|\widehat g(\xi)|\geq c_0\) 于 \(|\xi|\leq\delta\)。在区间
\[
 |x|\leq\delta\cdot w(\mu)^2
\]
上，\(\chi(x)=1\)，且 \(|S_\mu g(x)|=|\widehat g(-\mu x)|\geq c_0\)。于是
\[
 \|S_\mu g\|_2^2
       \geq2\delta c_0^2\cdot w(\mu)^2.
\]
这给出与参数无关的一维下界。

多维算子等于按各坐标先后应用 \(S_{\lambda\varepsilon_j}\)。对其他变量固定后使用一维上界，再用 Fubini 积分，可逐个得到乘积上界；先对紧支撑简单函数证明，再由密度延拓到 \(L^2\) 即可。取乘积单位输入 \(f(y)=\prod_jg(y_j)\)，输出也分解成 \(\prod_jS_{\lambda\varepsilon_j}g(x_j)\)，其平方范数由 Tonelli 分解为一维平方范数的乘积，得到下界。

因此每个坐标都在 \(\lambda\varepsilon_j\) 越过一时才开始贡献半个衰减幂次。零参数方向没有振荡增益；若所有有效参数都至少为一，才可将乘积改写为 \(\lambda^{-d/2}\prod_j\varepsilon_j^{-1/2}\)。
\end{solution}

% 自拟；反例与计算：混合矩阵处处可逆不保证全局单射，两个重复片仍满足局部算子衰减。
\Needspace{6\baselineskip}
\begin{exercise}\label{b2:ex:21-new-5}
在 \(\mathbb R_x^2\times\mathbb R_y^2\) 上令
\[
 F(s,t)=(e^s\cos t,e^s\sin t),
 \quad x=(s,t),\quad \Phi(x,y)=F(x)\cdot y.
\]
固定 \(0<\delta<\pi/4\)，取实函数 \(\chi_0\in C_c^\infty\bigl((-\delta,\delta)^2\bigr)\)，在原点附近恒为一；令 \(v=(0,2\pi)\)、\(\chi_1(x)=\chi_0(x-v)\)、\(\chi=\chi_0+\chi_1\)。再取非零 \(\eta\in C_c^\infty(\mathbb R^2)\)，对 \(\lambda\geq1\)、\(f\in L^2(\mathbb R^2)\) 定义
\[
 T_\lambda f(x)=\chi(x)\cdot
       \int e^{2\pi\mathrm i\lambda\Phi(x,y)}\cdot\eta(y)\cdot f(y)\,\mathrm dy.
\]
证明混合 Hesse 矩阵处处可逆，但 \(T_\lambda T_\lambda^*\) 的核不能在整个支撑上满足任何固定 \(N>0\) 的统一估计
\(C_N(1+\lambda|x-z|)^{-N}\)。另一方面，直接利用两个坐标片中的变量替换证明
\(\|T_\lambda\|_{2\to2}\leq C\lambda^{-1}\)。若 \(T_{\lambda,0}\) 表示仅用 \(\chi_0\) 的算子，进一步证明
\(\|T_\lambda f\|_2^2=2\|T_{\lambda,0}f\|_2^2\)。
\end{exercise}
\begin{solution}
直接求导给出
\[
 DF(s,t)=e^s
 \begin{pmatrix}\cos t&-\sin t\\\sin t&\cos t\end{pmatrix},
 \quad\det DF(s,t)=e^{2s}>0.
\]
按正文的约定，混合矩阵第 \((j,k)\) 项为 \(\partial_{y_j}\partial_{x_k}\Phi\)，此处它就是 \(DF(x)\)，所以处处可逆。但 \(F(x+v)=F(x)\)，映射并不全局单射。

相应的 \(TT^*\) 核是
\[
 K_\lambda(x,z)=\chi(x)\cdot\overline{\chi(z)}\cdot
       \int e^{2\pi\mathrm i\lambda(F(x)-F(z))\cdot y}
                                      \cdot|\eta(y)|^2\,\mathrm dy.
\]
取 \(x=0\)、\(z=v\)，两个截断值都为一，且相位差恒零，因此
\(K_\lambda(0,v)=\|\eta\|_2^2>0\)。而所提上界在这两点为 \(C_N(1+2\pi\lambda)^{-N}\to0\)，矛盾。

在第一个矩形中，极角限制于 \((-\delta,\delta)\)，所以 \(F\) 单射并为其像上的微分同胚。写 \(q=\eta\cdot f\)，则
\[
 T_{\lambda,0}f(x)=\chi_0(x)\cdot\widehat q(-\lambda F(x)).
\]
令 \(u=F(x)\)。该矩形中 \(e^{2s}\) 有正下界，且 \(\chi_0\) 有界，故
\[
 \begin{aligned}
 \|T_{\lambda,0}f\|_2^2
 &\leq C\int_{F\bigl((-\delta,\delta)^2\bigr)}
                      |\widehat q(-\lambda u)|^2\,\mathrm du\\
 &\leq C\lambda^{-2}\cdot\|q\|_2^2
 \leq C\lambda^{-2}\cdot\|\eta\|_\infty^2\cdot\|f\|_2^2.
 \end{aligned}
\]
输入 \(q\) 因紧支撑和 Cauchy--Schwarz 属于 \(L^1\cap L^2\)，所用 Fourier 表达合法。

两个截断支撑互不相交。在第二片作平移 \(x\mapsto x+v\)，有
\(\chi_1(x+v)=\chi_0(x)\) 和 \(F(x+v)=F(x)\)，所以第二片输出是第一片输出的平移副本。平方积分相加恰得到题中的二倍恒等式，继而得到整体 \(C\lambda^{-1}\) 界。失效的是跨越两片的单一梯度差下界；各片中的估计及其有限求和仍然有效。
\end{solution}

% 自拟；中等：不同输入输出维数中的秩模型，分离无振荡变量，计算准确范数并证明达到。
\Needspace{6\baselineskip}
\begin{exercise}\label{b2:ex:21-new-6}
设整数 \(r\geq1\)、\(m>r\)、\(n>r\)，实矩阵 \(B\in\mathbb R^{r\times r}\) 可逆。写
\(x=(x',x'')\in\mathbb R^r\times\mathbb R^{m-r}\)、
\(y=(y',y'')\in\mathbb R^r\times\mathbb R^{n-r}\)。给定非零
\(u\in L^2(\mathbb R^{m-r})\)、\(v\in L^2(\mathbb R^{n-r})\)，对 \(\lambda>0\) 及 \(f\in C_c^\infty(\mathbb R^n)\) 定义
\[
 T_\lambda f(x)
 =u(x'')\cdot\iint
       e^{2\pi\mathrm i\lambda x'\cdot By'}
       \cdot\overline{v(y'')}\cdot f(y',y'')\,\mathrm dy''\,\mathrm dy'.
\]
证明此算子唯一延拓为 \(L^2(\mathbb R^n)\to L^2(\mathbb R^m)\) 的有界算子，并求其准确范数。构造达到范数的非零 \(L^2\) 输入，说明衰减幂次为何由 \(r\) 决定，而不由 \(m\)、\(n\) 中任一个单独决定。
\end{exercise}
\begin{solution}
先对 \(y''\) 作配对，令
\[
 h(y')=\int\overline{v(y'')}\cdot f(y',y'')\,\mathrm dy''.
\]
对一般 \(L^2\) 输入，Fubini 与 Cauchy--Schwarz 保证此积分对几乎每个 \(y'\) 绝对收敛，并有
\[
 \|h\|_2\leq\|v\|_2\cdot\|f\|_2.
\]
在原来的光滑紧支撑输入上，\(h\) 还属于 \(L^1\)，原积分因而可识别为
\[
 T_\lambda f(x)=u(x'')\cdot\widehat h(-\lambda B^{\mathsf T}x').
\]
在 \(x'\) 上换元并使用 Plancherel，得到
\[
 \|T_\lambda f\|_2^2
 =\|u\|_2^2\cdot\lambda^{-r}\cdot|\det B|^{-1}\cdot\|h\|_2^2
 \leq\lambda^{-r}\cdot|\det B|^{-1}
             \cdot\|u\|_2^2\cdot\|v\|_2^2\cdot\|f\|_2^2.
\]
因此由 \(C_c^\infty\) 在 \(L^2\) 中的稠密性可唯一延拓；也可以直接用一般 \(L^2\) 的部分配对与 Plancherel 变换定义同一延拓。

取任意单位向量 \(h_0\in L^2(\mathbb R^r)\)，令
\[
 f(y',y'')=h_0(y')\cdot\frac{v(y'')}{\|v\|_2}.
\]
它是单位输入，部分配对恰为 \(h=\|v\|_2\cdot h_0\)，所以前面的不等式取等，得到
\[
 \boxed{\ \|T_\lambda\|_{2\to2}
 =\lambda^{-r/2}\cdot|\det B|^{-1/2}
                   \cdot\|u\|_2\cdot\|v\|_2\ }.
\]
相位 \(x'\cdot By'\) 的混合矩阵按输入变量作行指标，是一个左上块为 \(B^{\mathsf T}\)、其余为零的 \(n\times m\) 矩阵，秩为 \(r\)。变量 \(y''\) 只参与有界线性配对，变量 \(x''\) 只提供固定的 \(L^2\) 因子，两者都不增加振荡衰减。这个模型把混合相位的振荡方向与其余变量的配对作用明确分开。
\end{solution}

% 自拟；中等：坐标换元的酉归一化，推导两侧换元后的振幅及混合矩阵，检验局部估计的坐标一致性。
\Needspace{6\baselineskip}
\begin{exercise}\label{b2:ex:21-new-7}
设整数 \(d\geq1\)，\(U,V,X,Y\subset\mathbb R^d\) 为开集，\(F:U\to V\)、\(G:X\to Y\) 为光滑微分同胚，记
\(J_F=|\det DF|\)、\(J_G=|\det DG|\)。证明
\[
 (U_Ff)(x)=J_F(x)^{1/2}\cdot f\bigl(F(x)\bigr)
\]
给出 \(L^2(V)\to L^2(U)\) 的酉算子，并求其逆；\(U_G\) 同样定义。对 \(\lambda\geq1\) 及 \(f\in L^2(Y)\)，设
\[
 T_\lambda f(v)=\int_Y e^{2\pi\mathrm i\lambda\Phi(v,y)}
                                  \cdot a(v,y)\cdot f(y)\,\mathrm dy,
\]
其中实相位 \(\Phi\) 在 \(V\times Y\) 上光滑，\(a\in C_c^\infty(V\times Y)\)。求
\(\widetilde T_\lambda=U_FT_\lambda U_G^{-1}\) 的相位、振幅与混合 Hesse 矩阵，证明两算子的 \(L^2\) 范数完全相同，并说明局部非退化条件如何在换元下保持。不能只改变相位而省去振幅中的体积因子。
\end{exercise}
\begin{solution}
变量替换给出
\[
 \int_U J_F(x)\cdot|f\bigl(F(x)\bigr)|^2\,\mathrm dx
       =\int_V|f(v)|^2\,\mathrm dv.
\]
其逆公式为
\[
 (U_F^{-1}h)(v)
 =J_F\bigl(F^{-1}(v)\bigr)^{-1/2}\cdot h\bigl(F^{-1}(v)\bigr).
\]
代入并换元可验证逆映射也保范数，且两次复合为恒等，所以 \(U_F\) 为满射等距映射，即酉算子。\(G\) 的结论相同。

把逆公式代入原算子，再令 \(y=G(z)\)。输入端的逆平方根与 \(\mathrm dy=J_G(z)\,\mathrm dz\) 相乘，得到
\[
 \begin{aligned}
 \widetilde\Phi(x,z)&=\Phi\bigl(F(x),G(z)\bigr),\\
 \widetilde a(x,z)&=J_F(x)^{1/2}\cdot J_G(z)^{1/2}
                                \cdot a\bigl(F(x),G(z)\bigr),\\
 \widetilde T_\lambda h(x)
 &=\int_Xe^{2\pi\mathrm i\lambda\widetilde\Phi(x,z)}
                            \cdot\widetilde a(x,z)\cdot h(z)\,\mathrm dz.
 \end{aligned}
\]
原振幅的紧支撑经乘积微分同胚的逆像仍是紧集，故新振幅也是光滑紧支撑的。

令 \(M(v,y)\) 的第 \((j,k)\) 项为 \(\partial_{y_j}\partial_{v_k}\Phi\)，与正文采用同一约定。因两套换元分别只依赖自己的变量，链式法则没有额外的二阶交叉项，给出
\[
 \widetilde M(x,z)
 =DG(z)^{\mathsf T}\,M\bigl(F(x),G(z)\bigr)\,DF(x).
\]
两端矩阵可逆，因此 \(\widetilde M\) 可逆当且仅当对应的 \(M\) 可逆。固定紧支撑附近的局部估计仍可应用；变换后的具体导数界应按新坐标计算。

另一方面，酉映射把整个单位球双射到单位球，故直接有准确等式
\[
 \|\widetilde T_\lambda\|_{L^2(X)\to L^2(U)}
       =\|T_\lambda\|_{L^2(Y)\to L^2(V)}.
\]
这比仅分别应用两次估计更直接。平方根权重正是平方范数与变量替换相容所需的因子；例如在一维取实数 \(c\ne0\)，单独拉回 \(f\mapsto f(cx)\) 的范数为 \(|c|^{-1/2}\)，一般并不保范数。遗漏权重就改变了当前比较的算子。
\end{solution}

% 自拟；提高且依赖选读内容：对频率局部波包作Taylor误差估计，以实际能量集中检验半波传播正负号。
\Needspace{6\baselineskip}
\begin{optionalexercise}\label{b2:ex:21-new-8}
设整数 \(d\geq1\)，\(\omega\in\mathbb R^d\)、\(|\omega|=1\)，\(x_0\in\mathbb R^d\)，\(0<\alpha<1/2\)。取 \(h\in\mathcal S(\mathbb R^d)\)、\(\|h\|_2=1\)，并假定 \(\operatorname{supp}\widehat h\subset B(0,R)\)，其中 \(R>0\)。对 \(N\to\infty\)，定义单位波包
\[
 f_N(x)=N^{\alpha d/2}e^{2\pi\mathrm iN\omega\cdot x}
                       \cdot h\bigl(N^\alpha(x-x_0)\bigr).
\]
令 \(\sigma\in\{1,-1\}\)，半波算子 \(U_t^\sigma\) 的 Fourier 乘子为
\(e^{\sigma2\pi\mathrm it|\xi|}\)，其中 \(t\in\mathbb R\)。证明当 \(N^{1-\alpha}\geq2R\) 时，
\[
 \begin{aligned}
 v_N^\sigma(t,x)
 &=N^{\alpha d/2}e^{2\pi\mathrm i(N\omega\cdot x+\sigma tN)}
                    \cdot h\bigl(N^\alpha(x-x_0+\sigma t\omega)\bigr),\\
 \|U_t^\sigma f_N-v_N^\sigma(t,\cdot)\|_2
 &\leq C|t|N^{2\alpha-1}
                \bigl\|\,|\eta|^2\widehat h(\eta)\,\bigr\|_2,
 \end{aligned}
\]
其中 \(C\) 与 \(N,t\) 无关。由此证明，对固定 \(t\) 和任意有界连续函数 \(b\)，
\[
 \int b(x)\cdot|U_t^\sigma f_N(x)|^2\,\mathrm dx
       \longrightarrow b(x_0-\sigma t\omega).
\]
用所得位置说明正、负两个演化指数的传播方向；并指出允许时间依赖 \(N\) 时，上述近似保证收敛的时间尺度。
\end{optionalexercise}
\begin{solution}
调制、平移与缩放公式给出
\[
 \widehat f_N(\xi)
 =N^{-\alpha d/2}e^{-2\pi\mathrm ix_0\cdot(\xi-N\omega)}
                \cdot\widehat h\bigl(N^{-\alpha}(\xi-N\omega)\bigr).
\]
它只支撑于 \(\xi=N\omega+N^\alpha\eta\)、\(|\eta|\leq R\) 的区域。在题设范围内，连接 \(N\omega\) 与这些频率的线段距原点至少为 \(N/2\)。函数 \(\xi\mapsto|\xi|\) 的 Hesse 矩阵在该线段上的算子范数至多为 \(2/N\)，Taylor 积分余项因而给出
\[
 \bigl||N\omega+N^\alpha\eta|-N-N^\alpha\omega\cdot\eta\bigr|
       \leq C N^{2\alpha-1}|\eta|^2.
\]
将半波乘子中的长度换成线性项
\(N+\omega\cdot(\xi-N\omega)\)，Fourier 反演恰好得到题中的 \(v_N^\sigma\)：常数项产生 \(e^{\sigma2\pi\mathrm itN}\)，线性项使包络平移为 \(x-x_0+\sigma t\omega\)。

利用 \(|e^{\mathrm ia}-e^{\mathrm ib}|\leq|a-b|\)，再用 Plancherel 和变量替换 \(\xi=N\omega+N^\alpha\eta\)，得到
\[
 \|U_t^\sigma f_N-v_N^\sigma(t,\cdot)\|_2^2
 \leq C^2t^2N^{4\alpha-2}
       \int|\eta|^4\cdot|\widehat h(\eta)|^2\,\mathrm d\eta.
\]
开平方即为误差界。两种函数的 \(L^2\) 范数都是一，所以 Cauchy--Schwarz 还给出
\[
 \bigl\|\,|U_t^\sigma f_N|^2-|v_N^\sigma(t,\cdot)|^2\,\bigr\|_1
       \leq2\|U_t^\sigma f_N-v_N^\sigma(t,\cdot)\|_2.
\]
对固定 \(t\)，右侧因 \(2\alpha-1<0\) 而趋零。

包络的概率密度则可以直接换元：
\[
 \int b(x)\cdot|v_N^\sigma(t,x)|^2\,\mathrm dx
 =\int b\bigl(x_0-\sigma t\omega+N^{-\alpha}z\bigr)
                                      \cdot|h(z)|^2\,\mathrm dz.
\]
被积函数由 \(\|b\|_\infty\cdot|h|^2\) 控制，连续性与支配收敛给出极限 \(b(x_0-\sigma t\omega)\)。再用上面的密度 \(L^1\) 误差，即得真实演化的同一极限。

因此正指数 \(e^{2\pi\mathrm it|\xi|}\) 的包络中心为 \(x_0-t\omega\)，负指数的中心为 \(x_0+t\omega\)。若时间取 \(t=t_N\)，误差估计保证近似成立的充分条件是
\(|t_N|N^{2\alpha-1}\to0\)，即 \(|t_N|=o(N^{1-2\alpha})\)。这只是当前线性近似保证的尺度；在更长时间上，估计不能直接断言近似成立或失效。
\end{solution}
